\documentclass[11pt]{article}

\usepackage{maiacustom}

\begin{document}

\psettitle{Astronomy Question Bank}

These questions were carefully written or selected to prepare students for the astronomy team selection process in Brazil. Some questions are not original; those are marked with () before the statement. The question-bank template is the same one used by Professor Kevin Zhou. His work is invaluable, and many ideas in this set can be found in his handouts.

% \begin{psidea}{Idea title}{}
% Idea
% \end{psidea}

% \begin{psexample}{Example title}{}
% Example
% \end{psexample}

% \begin{pssolution*}{}{}
% Solution
% \end{pssolution*}

% \begin{psremark*}{Remark title}{}
% Remark
% \end{psremark*}

\section{Celestial Mechanics}
\pts{5}    
\begin{pproblem}
    Iaum and Sevla, two renowned physicists, intend to launch a space probe to explore the limits of Newtonian physics and the relativistic corrections required near a black hole. They need your help to study the \textit{effective potential} of such a body.
    \begin{enumerate}[label=\textbf{\alph*)}]
        \item The energy of a body of mass \(m\) orbiting a massive body of mass \(M\) can be written as:
             \[
             \frac{m\dot r ^2}{2} + V_{eff}(r)
             \]
             where \(\dot r\) is the radial velocity of the body. Find \(V_{eff}(r)\) in terms of \(G\), \(M\), \(L\), \(m\), and \(r\).

        \item According to Newtonian physics, what is the smallest radius at which a body can have a circular orbit around a black hole?
    
        \item What is the frequency of small radial oscillations, \(\omega_r\), about this radius?

        The gravitational potential near these bodies must include a relativistic correction and has the form:
        \[
        V(r) = -\frac{GMm}{r} - \frac{GML^2}{mc^2r^3}
        \]

        \item Explain why this correction allows a particle to fall into the center of a black hole, \(r=0\), and why this is impossible in Newtonian physics.
        
        \item What is the smallest value of \(L\) for which a particle can orbit the black hole on a circular orbit? What is the value of the radius in this case?
        
        \item Assuming \( \lim_{L \to \infty}\), find the smallest distance to which a particle on an open orbit can approach a black hole.  
    \end{enumerate}
    \begin{pssolution*}{}{}
        \begin{alternativas}
            \item In the classical form, we have:
            \[E = \frac{mv^2}{2}-\frac{GMm}{r}\]
            Note that we can write \(v^2 = v_\theta^2+ v_r^2\), where \(v_\theta\) and \(v_r\) are, respectively, the tangential and radial velocities of the body. Since \(v_r = \dot{r}\), we need an expression that relates \(v_\theta\) to another quantity. The best way to make this relation is by using the angular momentum, since:
            \[L = mrv_\theta \rightarrow v_\theta = \frac{L}{mr}\]
            Substituting into the energy expression, we have:
            \[E = \frac{m\dot{r}^2}{2}+\frac{L^2}{2mr^2} - \frac{GMm}{r}\]
            Since we want an answer of the form \(E = m\dot{r}^2/2 + V_{eff}(r)\), by analogy we have:
            \[\boxed{V_{eff}(r) = -\frac{GMm}{r}+\frac{L^2}{2mr^2}}\]

            \item The greatest limiting speed of a body is \(c\), disregarding relativistic effects:
            \[\frac{mc^2}{2}-\frac{GMm}{R} = -\frac{GMm}{2R}\]
            Isolating \(R\):
            \[\boxed{R = \frac{GM}{c^2}}\]

            \item For small oscillations, we have the approximation:
            \[\omega_r \approx \sqrt{\frac{V_{eff}''(R)}{m}}\]
            Calculating:
            \[V_{eff}'(R) = \frac{GMm}{r^2}-\frac{L^2}{mr^3}\]
            \[V_{eff}''(R) = -\frac{2GMm}{r^3}+\frac{3L^2}{mr^4}\]
            Using \(L = mvr\) (circular orbits) and \(v\equiv c\), we arrive at:
            \[V_{eff}''(R) = \frac{mc^6}{G^2M^2}\]
            With this, we can conclude that:
            \[\omega_r = \frac{c^3}{GM}\]

            \item When we take the limit \(\lim_{r\rightarrow 0}V_{eff}(r)\), in Newtonian physics the dominant term is \(L^2/2mr^2\), that is, an energy of \(+\infty\) would be required to reach the center of the black hole. However, after the relativistic correction, the dominant term is \(-GML^2/mc^2r^3\), that is, an energy of \(-\infty\) would be required to reach the center of the black hole. That is why it is not possible to reach the center of a black hole in Newtonian physics.
            
            \item On a circular orbit, \(r\) is constant, which implies that the force in the radial direction vanishes, that is:
            \[V_{eff}'(r) = 0\]
            Calculating:
            \[V_{eff}'(r) = \frac{GMm}{r^2}-\frac{L^2}{mr^3}+\frac{3GML^2}{mc^2r^4}=0\]
            Simplifying, we have:
            \[r^2-\frac{L^2}{GMm^2}r+\frac{3L^2}{m^2c^2}=0\]
            Note that the solution of this equation must be real, since \(r\) cannot be imaginary, which imposes the condition:
            \[\frac{L^4}{G^2M^2m^4} - \frac{12L^2}{m^2c^2} \geq 0\]
            In the equality case, the value of \(L\) that satisfies this equation is:
            \[\boxed{L = \frac{\sqrt{12}GMm}{c}}\]
            The value of \(r\) in this case is:
            \[r = \frac{L^2}{2GMm^2} = \boxed{r_{min} = \frac{6GM}{c^2}}\] 
            
            \item The complete expression for \(r\) is:
            \[r = \frac{\frac{L^2}{GMm^2}\pm\sqrt{\left(\frac{L^2}{GMm^2}\right)^2-\frac{12L^2}{m^2c^2}}}{2} = \frac{L^2-\sqrt{L^4-12(GMmL/c)^2}}{2GMm^2}\] 
            In the limit \(L\rightarrow \infty\), we have:
            \[r_{min} = \lim_{L \rightarrow \infty}\frac{L^2-\sqrt{L^4-12(GMmL/c)^2}}{2GMm^2} = \frac{6(GMm/c)^2}{2(GMm^2)} = \boxed{\frac{3GM}{c^2}} \] 
        \end{alternativas}
    \end{pssolution*}
\end{pproblem}

\pts{3}
\begin{pproblem} An important point in the study of orbits is what happens when the potential follows something unusual. The goal of this question is to study that phenomenon. Consider the angular momentum \(L\) and the mass of the body \(m\).
    \begin{alternativas}
        \item For a general orbit that has a potential \(V(r)\) of the form \(\psi r^k\), find the value \(r_0\) of a circular orbit for this potential.
        
        \item Find the frequency of small oscillations, \(\omega_r\), about this radius. 
        
        \item Let \(\omega_\theta\) be the angular frequency of the orbit. Find the value of the ratio \(\frac{\omega_r}{\omega_\theta}\). An orbit can be closed only if this ratio is an integer. Why? Find values of \(k\) for which the orbit is closed.
         
    \end{alternativas} 

    \begin{pssolution*}{}{}
        \begin{alternativas}
            \item Since on a circular orbit \(V_{eff}'(r) = 0\), we have:
            \[V_{eff}'(r) = -\frac{L^2}{mr^3}+k\psi r^{k-1} = 0\]
            \[k\psi r^{k-1} = \frac{L^2}{mr^3} \rightarrow \boxed{r_0 = \left(\frac{L^2}{k\psi m}\right)^{\frac{1}{k+2}}}\]
        
            \item Using \(\omega_r^2=V_{eff}''(r_0)/m\):
            \[V_{eff}''(r) = \frac{3L^2}{mr^4}+k(k-1)\psi r^{k-2}\]
            \[V_{eff}''(r) = r^{-4}\left(\frac{3L^2}{m}+k(k-1)\psi r^{k+2}\right)\]
            \[V_{eff}''(r) = r^{-4}\left(\frac{3L^2}{m}+k(k-1)\psi \frac{L^2}{k\psi m}\right)\]
            \[V_{eff}''(r) = \frac{L^2}{r_0^4m}(k+2)\]
            Thus, we have:
            \[\boxed{\omega_r = \frac{L}{r_0^2m}\sqrt{k+2}}\]

            \item By the definition of \(L\), we have:
            \[L = mr_0^2\omega_\theta \rightarrow \omega_\theta = \frac{L}{mr_0^2}\]
            Therefore:
            \[\frac{\omega_r}{\omega_\theta} = \sqrt{k+2}\]
            Note that, for the orbit to be a closed figure, this ratio must be an integer, so \(k+2\) must be a perfect square. Any value \(k = n^2-2 \ \forall \ n \in \mathbb{N}\) satisfies this equation.
        \end{alternativas}
    \end{pssolution*}
\end{pproblem}

\pts{4} %week 114
\begin{pproblem}
    The hodograph is one of the lesser-known ways of solving celestial-mechanics problems. Using this device, one can solve complex problems (involving, for example, finding orbital parameters, velocity parameters, and how to minimize the eccentricity of an orbit). Although the hodograph has many uses, the goal of this exercise is to prove the Hodograph Theorem, which states that, on closed orbits, the velocity forms a circle in vector space. We will prove this theorem in two ways.
    \begin{align*}
    \textbf{Method 1)}
    \end{align*}
    \begin{enumerate}[label=\textbf{\Roman*.}]
        \item Starting from Newton's second law and the conservation of angular momentum, find an expression for \(d\mathbf{v}/d\theta\). Leave your answer in terms of the angular momentum per unit mass, \(h\), and \(\mu\) = \(GM\). 
        \item Where is the center of this circle located? Taking the central body at the origin of the system, determine an expression for the distance between the central body and the center of the hodograph. 
    \end{enumerate}
    
    \begin{align*}
        \textbf{Method 2)}
    \end{align*}
    Method 2 uses more advanced notions of calculus and starts from the conservation of the \textit{Laplace-Runge-Lenz vector}, \(\mathbf{A}\).
    \begin{enumerate}[label=\textbf{\Roman*.}]
        \item The vector \(\mathbf{A}\) is given by the expression:
        \[
        \mathbf{A} = \mathbf{p}\times \mathbf{L} - GMm^2\mathbf{\hat{r}}
        \]
        The first step of the proof is to show that \(\mathbf{A}\) is constant. To do so, differentiate it with respect to time and draw your own conclusions.
        \item Since the vector \(\mathbf{L}\) is constant, \(\mathbf{A}\times\mathbf{L}\) is also a constant. Using the vector identity \(\mathbf{A}\times(\mathbf{B}\times\mathbf{C}) = \mathbf{B}(\mathbf{A}\cdot\mathbf{C}) - \mathbf{C}(\mathbf{A}\cdot\mathbf{B}) \), prove that \(\mathbf{v}\) forms a circle in vector space.    
    \end{enumerate}
    \begin{pssolution*}{}{}
        \textbf{Method 1)}

        \begin{enumerate}[label=\textbf{\Roman*.}]
            \item Newton's second law tells us that:
            \[-\frac{GMm}{r^2}\mathbf{\hat{r}} =  m\frac{d\mathbf{v}}{dt}\]
            By the definition of angular momentum:
            \[L = mr^2\frac{d\theta}{dt} \rightarrow \ \  \therefore dt = \frac{mr^2}{L}d\theta\]
            Substituting \(dt\) into Newton's second law:
            \[-\frac{GMm}{r^2}\mathbf{\hat{r}} = m\left(\frac{L}{mr^2}\right)\frac{d\mathbf{v}}{d\theta}\]

            Isolating \(dv/d\theta\), we have:

            \[\boxed{\frac{d\mathbf{v}}{d\theta} = -\frac{GMm}{L}\mathbf{\hat{r}} \equiv -\frac{\mu}{h}\mathbf{\hat{r}}}\]

            Note that this expression is a circle in vector space and has radius \(R = \mu/h\), since \(d\mathbf{v}/d\theta\) is constant and points in the radial direction.

            \item As we saw in the previous item, the hodograph forms a circle of radius \(R = \mu/h\). Use the following figure as support:
            \begin{figure}[H]
                \centering
                \includegraphics[width=0.4\linewidth]{imagens/hodografo1.png}
                \caption{Diagram of the hodograph}
            \end{figure}
            The point \(O\) represents the origin of the system (location of the central body). Note that the smallest distance between \(O\) and the hodograph must correspond to the smallest speed (apastron) and the largest distance must correspond to the largest speed (periastron). Also note that \(v_a + \overline{OC} = v_p - \overline{OC}\). Setting up the equations, we have:
            \[v_a + v_p = 2R \equiv \frac{2\mu}{h}\]
            \[v_p - v_a = 2\overline{OC}\]

            We can write \(v_p\) and \(v_a\) in terms of the angular momentum:
            \[L = ma(1-e)v_p \rightarrow v_p = \frac{h}{a(1-e)}\]
            \[L = ma(1+e)v_a \rightarrow v_a = \frac{h}{a(1+e)}\]
            
            Substituting into the first equation:
            \[\frac{h}{a}\left(\frac{1}{1+e}+\frac{1}{1-e}\right) = \frac{2\mu}{h}\]
            
            Here, we can isolate \(a\):
            \[\frac{h}{a}\left(\frac{2}{(1-e)^2}\right) = \frac{2\mu}{h}\]

            \[a = \frac{h^2}{\mu}\frac{1}{(1-e)^2}\]

            With this in mind, we can move on to the second expression:
            \[\frac{h}{a}\left(\frac{1}{1-e}-\frac{1}{1+e}\right) = 2\overline{OC}\]
            \[h\left(\frac{\mu(1-e)^2}{h^2}\right)\left(\frac{2e}{(1-e)^2}\right) = 2\overline{OC}\]
            
            Finally, solving for \(\overline{OC}\), we have:
            \[\boxed{\overline{OC} = \frac{\mu}{h}e}\]

            For an ellipse (\(e<1\)), the origin lies inside the circle; for a parabola (\(e=1\)), the origin lies on the boundary of the circle; and for a hyperbola (\(e>1\)), the origin lies outside the circle.
        \end{enumerate}

    
        \textbf{Method 2}
        \begin{enumerate}[label=\textbf{\Roman*.}]
            \item Differentiating \(\mathbf{A}\) with respect to time, we have:
            \[\frac{d\mathbf{A}}{dt} = \frac{d}{dt}(\mathbf{p}\times\mathbf{L}) - GMm^2\frac{d}{dt}\mathbf{\hat{r}}\] 
            When we differentiate a \textit{unit vector} with respect to time, we have:
            \[\frac{d\mathbf{\hat{r}}}{dt} = \mathbf{\omega}\times\mathbf{\hat{r}}\]
            where \(\mathbf{\omega}\) is the velocity with which the unit vector translates. Thus:
            \[\frac{d\mathbf{A}}{dt} = \left(\frac{d\mathbf{p}}{dt}\right)\times \mathbf{L} + \mathbf{p}\times\left(\frac{d\mathbf{L}}{dt}\right) - GMm^2(\mathbf{\omega}\times \mathbf{\hat{r}})\]
            Now we have a few considerations. First, note that \(d\mathbf{L}/dt=0\), since there are only central forces on the orbit. Second, \(d\mathbf{p}/dt = \mathbf{F}\). Furthermore, we have \(\mathbf{L} = mr^2\mathbf{\omega}\). Returning to the equation:

            \[\frac{d\mathbf{A}}{dt} = -\frac{GMm}{r^2}\mathbf{\hat{r}}\times(mr^2\mathbf{\omega}) - GMm^2(\mathbf{\omega} \times \mathbf{\hat{r}})\]
            \[\frac{d\mathbf{A}}{dt} = -GMm^2(\mathbf{\hat{r}}\times\mathbf{\omega})- GMm^2(\mathbf{\omega} \times \mathbf{\hat{r}})\]

            For any two vectors, \(\mathbf{i}\) and \(\mathbf{j}\), it holds that:
            \[i\times j = -j\times i\]

            Using this property, we can conclude that:

            \[\boxed{\frac{d\mathbf{A}}{dt} = 0}\]

            Therefore, \(\mathbf{A}\) is constant!

            \item Writing the expression and using \textit{BAC-CAB}, we have:
            
            \[\mathbf{A}\times\mathbf{L} = (\mathbf{p}\times\mathbf{L})\times\mathbf{L} - GMm^2\mathbf{\hat{r}}\times\mathbf{L}\]

            \[(\mathbf{A} +GMm^2\mathbf{\hat{r}})\times\mathbf{L} = - \mathbf{L}\times(\mathbf{p}\times\mathbf{L}) = \mathbf{p}(-\mathbf{L}\cdot\mathbf{L})-\mathbf{L}(-\mathbf{L}\cdot\mathbf{p})\]

            Since \(\mathbf{L}\) and \(\mathbf{p}\) are always perpendicular, \(\mathbf{L}\cdot\mathbf{p}=0\), so:

            \[(\mathbf{A}+GMm^2\mathbf{\hat{r}})\times\mathbf{L} = -mL^2\mathbf{v}\]

            Since both \(\mathbf{A}\) and \(\mathbf{\hat{r}}\) are perpendicular to \(\mathbf{L}\) and constant, \(\mathbf{A}+GMm^2\mathbf{\hat{r}}\) traces a circle. The cross product with \(\mathbf{L}\) merely rotates the circle by \(90^\circ\) and multiplies its radius by \(\mathbf{L}\). Therefore, \(\mathbf{v}\) traces a circle in vector space.
            
        \end{enumerate}
    \end{pssolution*}
\end{pproblem}

\pts{2} %problems of week 114
\begin{pproblem} One of the most fascinating phenomena of the Solar System is neutron stars. These are very small, extremely dense bodies that spin very fast. In this question, we will build a theoretical model for these stars.
    \begin{enumerate}[label=\textbf{\alph*)}]
        \item The Vela Pulsar, a neutron star located in the constellation Vela, has a frequency of \(11 \text{Hz}\), that is, it spins about its own axis \(11\) times per second, with equatorial radius \(R_e = 9.6\text{ km}\) and mass \(M = 1.88M_\odot\). Given this, calculate the ratio of its equatorial radius to its polar radius.
        \item What is the smallest rotation period the Vela Pulsar can have so that it does not break apart?
        \item Currently, the rate of change of the period of the Vela Pulsar is \(\frac{\Delta P}{\Delta t} = 1.25 \cdot 10^{-13}\text{s/s}\). This causes the pulsar to release an enormous amount of energy. The surface temperature of the pulsar is \(T \sim 10^6\text{ K}\). Calculate the order of magnitude of the ratio between the power emitted due to the increase of the period and the power emitted by the Stefan-Boltzmann law.
    \end{enumerate}
    \begin{pssolution*}{}{}
        \begin{alternativas}
            \item A star is always in hydrostatic equilibrium, so, writing the balance for a point on the equator and for one of the poles,
            \[-\frac{GM}{r_e}-\frac{1}{2}\omega^2r_e^2 = -\frac{GM}{r_p}\]

            Using \(\omega = 2\pi f\),

            \[\frac{GM}{r_p} = \frac{GM}{r_e}+2\pi^2f^2r_e^2\]

            \[\boxed{\frac{r_e}{r_p} = 1+\frac{2\pi^2f^2r_e^3}{GM} \approx 1 + 8.4\cdot10^{-6}}\]

            That is, a neutron star, even spinning absurdly fast, is still practically a perfect sphere!

            \item In the limiting situation in which it remains intact, we have:
            \[\frac{GM}{r_e^2} = \omega^2r_e\]
            \[\omega^2 = \frac{GM}{r_e^3} \rightarrow \boxed{T_{min} = \frac{2\pi r_e^3}{GM} = 2.2\cdot10^{-8}\text{ s}}\]

            \item The energy due to rotation has the form:
            \[E = \frac{1}{2}I\omega^2\]
            where \(I\) is the moment of inertia of the star (\(I = \frac{2}{5}Mr^2\)). Therefore, the dissipated power is:
            \[P = \frac{dE}{dt} = I\omega\frac{d\omega}{dt}\]

            Since \(\omega = 2\pi/T\),

            \[\frac{d\omega}{dT} = -\frac{2\pi}{T^2} \rightarrow d\omega = -\frac{2\pi}{T^2}dT\]
            
            The expression for the power as a function of the period becomes:

            \[P = -\frac{8\pi^2Mr_e^2}{5T^3}\frac{dT}{dt} = -9.1\cdot10^{29} \text{ W}\]

            Note that, by conservation of energy, all of this power must be converted into luminosity, so \(L_{rot} = -P = 9.1\cdot10^{29}\text{ W}\). The power coming from thermal radiation is given by:

            \[L_{rad} = 4\pi R^2\sigma T^4 = 6.56 \cdot 10^{25}\text{ W}\]

            With this, we can conclude that \(99.99\%\) of its luminosity comes from the rotation.
        \end{alternativas}
    \end{pssolution*}    
\end{pproblem}

\pts{3}
\begin{pproblem} (Adapted from PPP) Geométrio, Paulinho, and Hirata love chocolate. One day, they were exploring a galaxy when they came across a giant ball of chocolate. The three quickly began to eat the ball, first making a straight line from the point \(P\) to the center of the ball at \(O\) (diagram \(1\)). After that, Hirata falls, colliding at the point \(O\), without braking along the way. Surprisingly, he remains alive and is rescued by Geométrio and Paulinho. After that, still hungry, they continue to eat the giant chocolate sphere and leave a spherical hole of diameter \(\overline{PO}\) (diagram \(2\)). Hirata is very clumsy and ends up falling again, starting from the point \(P\) and going to \(O\) once more.
    \begin{enumerate}[label=\textbf{\alph*)}]
        \item Find the ratio of Hirata's impact speeds in cases \(1\) and \(2\).
        \item Find the ratio of Hirata's fall times in cases \(1\) and \(2\).
    \end{enumerate}
    \begin{paracol}{2}
        \begin{figure}[h]
            \centering
            \includegraphics[width=0.31\textwidth]{imagens/q5(1).png}
            \caption{Diagram \(1\)}
        \end{figure}
        \switchcolumn
        \begin{figure}[h]
            \centering
            \includegraphics[width=0.3\textwidth]{imagens/q5(2).png}
            \caption{Diagram \(2\)}
        \end{figure}
    \end{paracol} 

    \begin{pssolution*}{}{}
        \begin{alternativas}
            \item The gravity inside a planet is given by:
            \[\oint\mathbf{g}\cdot d\mathbf{S} = -4\pi GM_{int}\]
            \[4\pi r^2g = -4\pi G \rho \frac{4\pi}{3}r^3\]

            That is, for case \(1\), \(g \propto -\rho r \rightarrow g = -k\rho r\). Using Newton's second law:
            \[F = mv\frac{dv}{dr} = -mkr\]
            Integrating:
            \[\int_0^{v_1} v'dv' =-k\rho\int_0^R rdr\]

            \[\frac{v_1^2}{2} = k\rho\frac{R^2}{2} \rightarrow \boxed{v_1 = R \sqrt{k\rho}}\]

            Let \(\mathbf{j}\) be the vector from the center of the planet to the center of the cavity. Using superposition:
            \[a = -k\rho\mathbf{r} - k(-\rho)(\mathbf{r}-\mathbf{j}) = -k\rho\mathbf{j}\]
            That is, in case \(2\), the acceleration is constant and has magnitude \(|a| = \frac{k\rho R}{2}\).
            
            Using Torricelli:
            \[v_2^2 = 2aR k\rho R^2 \rightarrow \boxed{v_2 = R\sqrt{k\rho}}\]
            Therefore, \(v_1/v_2 = 1\).

            \item In the first case, note that we have simple harmonic motion, since the force is proportional to \(r\). The time to reach the bottom equals \(1/4\) of the period:
            \[\ddot{r} = -k\rho r \rightarrow \omega^2 = k\rho \rightarrow T = \frac{2\pi}{\sqrt{k\rho}} \rightarrow t_1 = \frac{T}{4} = \frac{\pi}{2\sqrt{k\rho}}\]

            In the second case, we have a constant acceleration, so:
            \[t_2 = \frac{v_2}{a_2} = \frac{2R\sqrt{k\rho}}{k\rho R} = \frac{2}{\sqrt{k\rho}}\]

            Therefore, the ratio \(t_1/t_2\) is given by:
            \[\boxed{\frac{t_1}{t_2} = \frac{\pi}{4}}\]

        \end{alternativas}
    \end{pssolution*}
\end{pproblem}

\pts{5}
\begin{pproblem} Calculating escape speeds in certain situations can be more complicated than it seems. One example is calculating the launch speed a rocket needs in order to reach a given location.
    \\
    \textbf{Hint:} The following problems require you to think carefully about the most appropriate reference frame. Do not underestimate the question.
    \\
    \textbf{Data:} A body can orbit the surface of the Earth with speed \(v_0 = 7.9\text{ km/s}\), the orbital speed of the Earth around the Sun is \(u_0 = 29.7\text{ km/s}\), and the radius of the Earth is negligible compared with the Earth-Sun distance. Also assume that, upon leaving the Earth's gravitational field, the distance between the probe and the Sun is the same as the distance between the Earth and the Sun.
    \begin{enumerate}[label=\textbf{\alph*)}]
        \item What is the smallest launch speed a rocket needs in order to reach the Sun, assuming that it applies only one impulse? (To check, \(v = 31.8\) km/s)
        \item What is the smallest launch speed a rocket needs in order to escape the Solar System? (To check, \(v = 16.7\) km/s)
        \item How does the answer to item \textbf{a)} change if the rocket can apply a second, very small impulse at some point of its orbit?
    \end{enumerate}

    \begin{pssolution*}{}{}
        \begin{alternativas}
            \item For this to happen, after leaving the Earth, in the Sun's frame the probe must be at rest and, consequently, in the Earth's frame it must have velocity \(-u_0\). Using conservation of energy:
            \[\frac{mv^2}{2}-\frac{GM_\oplus m}{R_\oplus} = \frac{mu_0^2}{2}\]
            \[v^2 = u_0^2+\frac{2GM_\oplus}{R_\oplus}\]

            Note that the speed \(v_0\) can also be obtained by equating the centripetal acceleration to the gravitational acceleration:
            \[\frac{mv_0^2}{R_\oplus} = \frac{GM_\oplus m}{R_\oplus^2} \rightarrow v_0^2 = \frac{GM_\oplus}{R_\oplus}\]

            Substituting,

            \[v^2=u_0^2+2v_0^2 \rightarrow \boxed{v = 31.8\text{ km/s}}\]

            \item In this case, upon leaving the Earth's gravitational field, in the Sun's frame the body must have speed \(\sqrt{2}u_0\), so in the Earth's frame it must have \((\sqrt{2}-1)u_0\). Again, by conservation of energy:
            
            \[\frac{mv^2}{2}-\frac{GM_\oplus m}{R_\oplus} = \frac{m(\sqrt{2}-1)^2u_0^2}{2}\]

            \[v^2 = (\sqrt{2}-1)^2u_0^2+2v_0^2 \rightarrow \boxed{v = 16.7 \text{ km/s}}\]

            \item Imagine that we send the probe to infinity, that is, we launch it with speed \(v=16.7\text{ km/s}\), and after that we apply an infinitesimal impulse toward the Sun. This makes the probe travel all the way back and, consequently, reach the Sun.
        \end{alternativas}
    \end{pssolution*}
\end{pproblem}

\pts{3} % problems of week 115
\begin{pproblem} (Kevin Zhou) The rocket equation is given by:
    \[
    v = u\ln\left(\frac{M_0}{M}\right)
    \]
    where \(v\) is the speed of the rocket, \(u\) the relative speed at which the rocket ejects fuel, \(M_0\) the initial mass of the rocket, and \(M\) its current mass.
    \begin{enumerate}[label=\textbf{\alph*)}]
        \item Derive the rocket equation starting from conservation of momentum.
        \item Assuming \(u\) is fixed throughout the rocket's trip, what must the value of \(u\) be so that the rocket goes from \(0\) to \(v\) using the least possible fuel? (You will need to solve numerically for \(v/u\))
        \item How must the rocket equation be corrected for a region of space with a constant gravitational field, \(\mathbf{g}\), pointing opposite the rocket's motion? Take \(\eta = dm/dt = \text{constant}\).
    \end{enumerate}

    \begin{pssolution*}{}{}
        \begin{alternativas}
            \item Consider that the rocket has mass \(M\) and velocity \(v\). It ejects fuel at a relative speed \(u\). Let \(dm\) be the mass of ejected gas:
            \[dp = (v-u)dm\]
            But, by definition,
            \[dp = vdm + mdv\]
            Combining these equations:
            \[mdv = -udm\]

            Integrating,
            \[\ln m = -\frac{v}{u} + C\]

            We know that when \(m = M_0\), we have \(v = 0\), so \(C = \ln M_0\). Thus:

            \[v = u \ln\left(\frac{M_0}{M}\right)\]

            As expected.

            \item The kinetic energy of the gas is given by:
            
            \[K_{g} =\frac{1}{2}M_gu^2 = \frac{(M_0-M)u^2}{2}\]

            We can find an expression for \(M\) using the rocket equation:

            \[\frac{M_0}{M} = e^{v/u} \rightarrow M_0 = Me^{v/u}\]

            Returning to the previous formula:
            
            \[K_g = \frac{M(e^{v/u}-1)u^2}{2}\]

            Here, \(M\) is equivalent to the ``structural'' mass of the rocket.

            To minimize the amount of energy spent on fuel, we have:

            \[\frac{dK_g}{du}=0  \rightarrow -e^{v/u}v +2u(e^{v/u}-1)=0\]

            Defining \(x=v/u\),

            \[e^xx = 2(e^x-1) \rightarrow x = 2(1-e^{-x})\]

            Solving by iteration, we find:

            \[x \approx 1.59 \rightarrow \boxed{u \approx \frac{v}{1.59}} \]

            \item In this case, we have:
            
            \[dp = (v-u)dm - mgdt = mdv +vdm\]
            \[\frac{dm}{m} = -\frac{dv}{u} - \frac{g}{u}dt\]

            Integrating,

            \[\ln\left(\frac{M}{M_0}\right) = -\frac{v}{u} - \frac{gt}{u}\]

            Since \(\eta\) is constant, \(M = M_0 -\eta t \rightarrow t = (M_0-M)/\eta\).

            \[\ln\left(\frac{M}{M_0}\right) = -\frac{v}{u} - \frac{g}{u\eta}(M_0-M)\]

            Solving for \(v\),

            \[\boxed{v = u\ln\left(\frac{M_0}{M}\right)-\frac{g}{\eta}(M_0-M)}\]
        \end{alternativas}
    \end{pssolution*}
\end{pproblem}  

\pts{4}
\begin{pproblem} One of the most common rocket propellants is a mixture of liquid hydrogen and liquid oxygen. When it begins to burn, the following chemical reaction occurs:
    \[
    2H_2 + O_2 \rightarrow 2H_2O
    \] 
    For each mole of hydrogen, this reaction releases \(241.8 \text{ kJ}\) of energy. Throughout the question, assume that all of this energy is used to move the rocket.
    \begin{enumerate}[label=\textbf{\alph*)}]
        \item A space expedition is to be carried out in such a way that a Hohmann transfer is required to launch a rocket from Earth to Mars. Calculate the total velocity change \(\Delta v\) needed to perform this maneuver.
        \item From the \(\Delta v\) calculated above, estimate how many tons of propellant must be used to perform this maneuver for a rocket that ejects propellant at speed \(u = 3.0\) km/s and has structural mass \(M = 140\cdot10^3\) kg.
    \end{enumerate}
    \begin{pssolution*}{}{}
        \begin{alternativas}
            \item For the first impulse, we have:
            \[v_{0,1}^2=\frac{GM}{R_T}\]
            \[v_{f,1}^2 = GM\left(\frac{2}{R_T}-\frac{1}{a}\right)\]
            The transfer orbit has \(a=(R_T+R_M)/2\),
            \[v_{f,1}^2 = \frac{2GM}{R_T}\left(\frac{R_M}{R_T+R_M}\right)\]

            For the second impulse,
            \[v_{0,2}^2 = GM\left(\frac{2}{R_M}-\frac{1}{a}\right) = \frac{2GM}{R_M}\left(\frac{R_T}{R_T+R_M}\right)\]

            \[v_{f,2}^2 = \frac{GM}{R_M}\]

            With \(\Delta v = \Delta v_1 + \Delta v_2 = (v_{f,1}-v_{0,1})+(v_{f,2}-v_{0,1})\), we can write:

            \[\boxed{\Delta v = \sqrt{\frac{2GM}{R_T}\left(\frac{R_M}{R_T+R_M}\right)}-\sqrt{\frac{GM}{R_T}} + \sqrt{\frac{GM}{R_M}} - \sqrt{\frac{2GM}{R_M}\left(\frac{R_T}{R_T+R_M}\right)}}\]
            
            Using \(M = M_\odot = 2\cdot10^{30}\) kg, \(R_M = 1.5\) AU, and \(R_T = 1\) AU, we find the numerical value \(\Delta v \approx 5.42 \) km/s.

            \item Using the rocket equation, we have:
            
            \[v = u\ln\left(\frac{M_0}{M}\right) \rightarrow \frac{M_0}{M}  = 1+ \frac{\mu}{M}= e^{v/u}\]

            where \(\mu\) is the mass of gas.

            \[\boxed{\mu = M(e^{v/u}-1) \approx 7.12 \cdot 10^5 \text{ kg}}\]

            This is consistent with reality, since the fuel mass is \(\approx 83\%\) of the total mass of the rocket.
        \end{alternativas}
    \end{pssolution*}
\end{pproblem}

\pts{3} % problems of week 115
\begin{pproblem} Marisso was studying a binary system with inclination \(i\). He found that the largest redshift coming from star \(1\) was \(z_1\). Given this, find an expression for the mass of star \(2\), \(m_2\). Leave your answer in terms of \(z_1\), the binary period \(P\), the inclination \(i\), the mass ratio \(\lambda  = m_2/m_1\), and fundamental constants. Assume circular orbits.
    \begin{pssolution*}{}{}
        The speed at which \(1\) orbits the CM is:

        \[v_1 = \frac{2\pi a_1}{P}\]

        By geometry, the radial velocity is \(v_{r,1}=v_1\sin i\), so:

        \[z_1 \approx \frac{v_{r,1}}{c} = \frac{2\pi a_1\sin i}{Pc}\]

        From the center-of-mass theorem:

        \[m_1a_1=m_2a_2 \rightarrow a_2 = \frac{a_1}{\lambda}\]

        Using Kepler's third law:

        \[\frac{(a_1+a_2)^3}{P^2} = \frac{G(m_1+m_2)}{4\pi^2}\]

        \[\frac{a_1^3(1+1/\lambda)^3}{P^2} = \frac{Gm_2(1+1/\lambda)}{4\pi^2}\]

        Isolating \(m_2\):

        \[m_2 = \frac{4\pi^2a_1^3(1+1/\lambda)^2}{GP^2}\]

        \(a_1\) can be obtained by evaluating the redshift formula, \(a_1 = Pcz_1/2\pi\sin i\):

        \[m_2 = \frac{4\pi^2}{GP^2}\left(\frac{z_1Pc}{2\pi\sin i}\right)^3(1+1/\lambda)^2\]

        Finally, simplifying:

        \[\boxed{m_2 = \frac{z_1^3c^3P}{2\pi G \sin^3i}(1+1/\lambda)^2}\]
    \end{pssolution*}
\end{pproblem}

\pts{2}
\begin{pproblem} What is the ratio between \(a)\) the gravitational forces caused by the Sun and by the Moon on the surface of the Earth? And \(b)\) the tidal forces caused by the same?
    \begin{pssolution*}{}{}
        \begin{alternativas}
            \item The gravitational force has the form:
            \[F_G = \frac{GMm}{d^2}\]

            Thus,

            \[\frac{F_\odot}{F_L} = \frac{M_\odot}{M_L}\left(\frac{d_{T-L}}{d_{T-\odot}}\right)^2 \boxed{\approx 178.33}\]

            \item For the tidal forces, we have:
            
            \[F_M \propto \frac{m}{d^3}\]

            \[\therefore \frac{F_\odot}{F_L} = \frac{M_\odot}{M_L}\left(\frac{d_{T-L}}{d_{T-\odot}}\right)^3 \boxed{\approx 0.456}\]

        \end{alternativas}
    \end{pssolution*}
\end{pproblem}

\pts{4}
\begin{pproblem} (Morin 7.7) A particle of mass \(m\) travels on a hyperbolic orbit with a mass \(M\) fixed at one of the foci. The speed at infinity is \(v_0\) and the impact parameter is \(b\).
    \begin{alternativas}
    \item Show that the deflection angle of the particle is given by:
    \[
    \phi = \pi-2\tan^{-1}(\gamma b)
    \]
    where \(\gamma \equiv v_0^2/GM\).
    \item Let \(d\sigma\) be the cross section of the particle (measured when it is at infinity) that is deflected into a solid angle of size \(d\Omega\). Show that:
    \[
    \frac{d\sigma}{d\Omega} = \frac{1}{4\gamma^2\sin^4(\phi/2)}
    \] 
    As a point of curiosity, this quantity is called the \textit{differential cross section}.
    \end{alternativas}

    \begin{pssolution*}{}{}
        \begin{alternativas}
            \item The diagram of the situation is the following:
            
            \begin{figure}[H]
                \centering
                \includegraphics[width=0.87\linewidth]{imagens/orbita hiperbolica esquema 1.png}
                \caption{Diagram of the hyperbolic orbit}
                \label{fig.1}
            \end{figure}

            To find the deflection angle \(\phi\), we will use the impulse theorem, considering the horizontal and vertical components of the attractive force:

            \[F_x = -\frac{GMm}{r^2}\cos\theta, \ \ F_y = -\frac{GMm}{r^2}\sin\theta\]

            The impulse theorem tells us that:
            \[\int F_i dt = \Delta p\]

            where \(p_i\) is the momentum of the particle in the direction \(i\).

            On the \(x\)-axis, we have:

            \[-\int \frac{GMm}{r^2}\cos\theta dt = \Delta p_x \]

            Note that integrating with respect to time makes things inconvenient, because finding how \(\theta\) and \(r\) vary in time is complicated. For this reason, we turn to the definition of angular momentum. Using \(L = mr^2d\theta/dt\), we have:

            \[dt = \frac{mr^2}{L}d\theta\]

            Substituting this value into our integral:
            \[-\int \frac{GMm^2}{L}\cos\theta d\theta = \Delta p_x\]

            Notice that our integral goes from \(\theta = 0\), the situation at infinity, to \(\theta = \pi - \phi\), the situation at infinity after the deflection. Thus:

            \[-\frac{GMm^2}{L}\int_0^{\pi-\phi}\cos\theta d\theta = -\frac{GMm^2}{L}\sin\theta\big|_0^{\pi-\phi} = -\frac{GMm^2}{L}\sin(\pi-\phi) = \Delta p_x\]

            Analogously for \(y\):

            \[-\frac{GMm^2}{L}\int_0^{\pi-\phi}\sin\theta d\theta = \Delta p_y\]
            \[=\frac{GMm^2}{L}\cos\theta\big|_0^{\pi-\phi} \equiv -\frac{GMm^2}{L}(1-\cos(\pi-\phi)) = \Delta p_y\]
        
            Also note that \(\Delta p_x = m(v_{f,x}-v_0), \ \Delta p_y = m(v_{f,y}-0)\). We can also find a relation between \(\phi\) and the velocities and, by pure geometry, we obtain:
            \[\tan\phi =-\frac{v_{f,y}}{v_{f,x}}\]
            Since there are only central forces, \(L\) is constant, and from the initial situation it equals \(L = mbv_0\). Now, turning to the calculation:

            \[-\frac{GM}{bv_0}\sin(\pi-\phi) = v_{f,x}-v_0, \ \ -\frac{GM}{bv_0}(1-\cos(\pi-\phi))=v_y\]
            
            With these two equations, we obtain:

            \[\frac{v_{f,y}}{v_{f_x}}= -\frac{GM(1-\cos(\pi-\phi))}{bv_0^2-GM\sin(\pi-\phi)}\]

            That is:

            \[\tan\phi = \frac{GM(1-\cos(\pi-\phi))}{bv_0^2-GM\sin(\pi-\phi)}\]

            Solving for \(\phi\), we obtain the expected result. Another way to carry out the same question is to think purely about the geometry of the situation:
            \begin{figure}[H]
                \centering
                \includegraphics[width=0.7\linewidth]{imagens/morin hiperbole.png}
                \caption{Source: Morin 7.7}
            \end{figure}

            Here, it is evident that:

            \[\phi = \pi - 2\tan^{-1}\left(\frac{b}{a}\right)\]

            where \(b\) is the impact parameter and \(a\) the semiaxis of the hyperbola. From the theory of conics, we have:
            \[\frac{b}{a} = \sqrt{e^2-1} = \sqrt\frac{2EL^2}{m(GMm)^2} = \frac{v_0^2b}{GM} \equiv \gamma b\]

            In both approaches, we arrive at the same result:

            \[\boxed{\phi = \pi - \tan^{-1}(\gamma b)}\]

            as desired.

            \item Consider a ring of thickness \(db\) and radius \(b\). Now consider a very large sphere, centered at \(M\). Any particle that passes through the cross section of the ring of radius \(b\) will hit this sphere making an angle \(\phi\) with the \(x\)-axis, with an angular separation \(d\phi\). Using \(d\cot x/dx = -1/\sin^2x\), we have:
            
            \[\big|\frac{db}{d\phi}\big| = \frac{1}{2\gamma \sin^2(\phi/2)}\]

            The cross-sectional area is given by \(d\sigma = 2\pi bdb\) and the solid angle on a sphere is given by \(d\Omega = 2\pi \sin\phi d\phi\). Therefore:

            \[\frac{d\sigma}{d\Omega} = \frac{2\pi b}{2\pi \sin \phi}\frac{db}{d\phi} = \left(\frac{(1/\gamma)\cot(\phi/2)}{2\sin(\phi/2)\cos(\phi/2)}\right)\left(\frac{1}{2\gamma\sin^2(\phi/2)}\right)\]

            \[\boxed{\frac{d\sigma}{d\Omega} = \frac{1}{4\gamma^2\sin^4(\phi/2)}}\]
        \end{alternativas}            
    \end{pssolution*}
\end{pproblem}

\begin{psidea}{}{}
    In certain questions, the use of vectors is appropriate. One example is the question below. To make things easier, when finding the inclination of the orbit, look at the following figure:

    \begin{figure}[H]
        \centering
        \includegraphics[width=0.7\linewidth]{imagens/vetores ecliptica.png}
        \caption{Vectors}
    \end{figure}

    Here, \(C\) represents a body, \(T\) the Earth, and \(S\) the Sun. The angles \(\lambda\) and \(b\) denote, respectively, the ecliptic longitude and the ecliptic latitude. \(\vec{r}\) is the vector between the Sun and the body, \(\vec{r}_T\) the vector between the Sun and the Earth, and \(\vec{d}\) the vector between the Earth and the body.
    From the figure,

    \[\vec{r} = r(\cos b\cos\lambda \hat{x}+ \cos b \sin\lambda \hat{y} + \sin b \hat{z})\]

    \[\vec{r}_T = r_T(\cos(\omega_\oplus t)\hat{x}+\sin(\omega_\oplus t)\hat{y})\]

    Thus, \(\vec{d}=\vec{r}-\vec{r}_T\):

    \[
    \vec{d} = \big(r \cos b \cos \lambda - r_T \cos (\omega_\oplus t)\big) \hat{x} 
    + \big(r \cos b \sin \lambda - r_T \sin (\omega_\oplus t)\big) \hat{y} 
    + r \sin b \, \hat{z}
    \]

    Calculating the magnitude of the vector \(\vec{d}\), we obtain:

    \[
    d = |\vec{d}| = \sqrt{r_T^2 + r^2 - 2 r_T r \cos b \cos (\lambda - \omega_\oplus t)}
    \]

    Therefore:

    \[
    \boxed{\cos b \cos (\lambda - \omega_\oplus t) = \frac{r_T^2 + r^2 - d^2}{2 r_T r}}
    \]
\end{psidea}

\pts{5}
\begin{pproblem}(Problem set \(1\) - Vinhedo 2024) 
    CR4b-2023 is a probe on a heliocentric orbit that was built and launched in 2023 by the young prodigy Caranguejo. The goal of the probe was to study solar storms and collect data for Caranguejo to analyze at his observatory in Espírito Santo. CR4b-2023, however, during one of its expeditions toward the Sun, was hit by a solar storm and suffered serious damage. Because of this, the probe became disoriented and thus had all of its orbital parameters changed, so that Caranguejo no longer knew its location.

    Aiming to track the position of CR4b-2023 again, Caranguejo used his observatory to collect the values of the angular separations \(\Delta\phi\) between the probe and the Sun and the values of the angular diameter \(\theta_S\) of the probe, all as a function of time \(t\). The table obtained by Caranguejo is shown below.

    \begin{table}[H]
        \centering
        \begin{tabular}{|c|c|c|}
            \hline
            \(\Delta\phi\) (Degrees) & \(\theta_S\) (mas) & \(t\) (Days) \\
            \hline
            0.000 & 99.27 & 0.00 \\
            4.889 & 103.67 & 1.79 \\
            17.354 & 104.72 & 7.43 \\
            29.015 & 93.06 & 20.62 \\
            27.793 & 81.18 & 25.41 \\
            13.460 & 79.77 & 45.98 \\
            11.539 & 77.97 & 51.61 \\
            2.466 & 76.44 & 54.53 \\
            2.656 & 74.84 & 55.31 \\
            7.177 & 73.90 & 57.19 \\
            10.800 & 70.63 & 58.43 \\
            22.464 & 72.92 & 91.67 \\
            13.823 & 80.50 & 130.48 \\
            \hline
        \end{tabular}
        \caption{Values measured by Caranguejo}
    \end{table}

    Assume that, at the initial instant \(t_0 = 0\) when the probe is hit by the solar storm, the Sun was exactly at the point of Libra and the motion of the probe was ascending relative to the ecliptic plane. Also assume that the radius of the probe, supposed spherical, is given by \(R = 30\) km. For this question, an error analysis is not required (although it is important to try to use methods aimed at reducing statistical errors). Based on what was presented:

    \begin{enumerate}[label=(\alph*)]
        \item Calculate the orbital parameters of the orbit of CR4b-2023, that is, its semimajor axis \(a\), eccentricity \(e\), inclination \(i\), longitude of the ascending node \(\Omega\), and argument of perihelion \(\omega\). As in all data-analysis questions, you must provide clearly labeled data tables, clearly labeled graphs, and enough formula derivations to make clear what you measured and how you are deriving your results in order to reduce statistical errors.
        
        \item Taking \((x', y', z')\) as the coordinates of CR4b-2023 in a right-handed Cartesian system in which the \(x'y'\) plane lies in the plane of the probe's orbit, the Sun is at the origin, and the \(x'\) axis points toward the vernal point, write, in a table, the values of at least 5 distinct points \((x', y', z')\). Based on the data found, sketch, on a graph, the orbit of CR4b-2023, indicating the coordinates of its center.
        
        \item Now considering a new right-handed coordinate system \((x, y, z)\), in which the Sun is at the origin, the \(xy\) plane represents the ecliptic plane, and the \(x\) axis points toward the vernal point, write the coordinates of the center of the orbit of CR4b-2023. Finally, calculate the distance \(d_{TC}\) between the center of Earth's orbit and the center of the probe's orbit.
    \end{enumerate}

    \begin{pssolution*}{}{}
        \begin{alternativas}
            \item The idea of this part of the question is to find a relation between the days and the distance from the probe to the Sun. Using basic trigonometry, we can say that the distance from the probe to the Earth is given by:
            \[d_{S-\oplus} = R/\tan^{-1}\left(\frac{\theta_S}{2}\right)\]
            where \(R\) is the radius of the probe.
            Here, we can use the law of cosines to relate this distance to the distance from the probe to the Sun. Considering the Earth's orbit to be circular with radius \(d_{\odot-\oplus}\) and defining \(d\) as the distance between the probe and the Sun, we have:
            \[d^2 = d_{\odot-\oplus}^2+d_{S-\oplus}^2-2d_{\odot-\oplus}d_{S-\oplus}\cos\Delta\phi\]      
            
            With this, we can build the following table, with the values of \(d\) in AU and the values of \(t\) in days.

            \[
                \begin{array}{|c|c|}
                \hline
                \text{t (Days)} & \text{d (UA)} \\
                \hline
                0.00 &    0.17 \\
                1.79 &    0.22 \\
                7.43 &    0.34 \\
                20.62 &    0.50 \\
                25.41 &    0.49 \\
                45.98 &    0.24 \\
                51.61 &    0.22 \\
                54.53 &    0.09 \\
                55.31 &    0.12 \\
                57.19 &    0.18 \\
                58.43 &    0.27 \\
                91.67 &    0.44 \\
                130.48 &    0.25 \\
                \hline
                \end{array}
            \]

            Plotting a graph of \(t\) versus \(d\), we have:
            
            \begin{figure}[H]
                \centering
                \includegraphics[width=0.85\linewidth]{imagens/graficoorbita.png}
            \end{figure}

            Now let us get to work and find the requested parameters: semimajor axis, \(a\), eccentricity, \(e\), inclination, \(i\), longitude of the ascending node \(\Omega\), and argument of perihelion, \(\omega\). The first two can be obtained directly from the data we have analyzed so far.

            Analyzing the data, we can see that the smallest distance between the probe and the Sun is given by \(a_p = 0.09\) AU and the largest distance is given by \(a_a = 0.49\) AU. These distances must correspond to perihelion and aphelion, respectively. Using the definition:
            \[a_p = a(1-e), \ a_a=a(1+e) \rightarrow \frac{a_p}{a_a}=\frac{1-e}{1+e} \approx 0.18\]
            
            Solving, we obtain \(e \approx 0.67\).

            Using \(a = \frac{a_p+a_a}{2}\), we obtain \(a\approx 0.30\) AU.

            To find the angular factors, we have to think a bit more. Starting with \(\Omega\), note that at the initial instant, at \(t=0\), the angular separation between the Sun and the probe is \(\Delta \phi = 0\). The statement tells us that at this moment the Sun is at the point of Libra, so the probe is also at the point of Libra and, since its motion was ascending, we have \(\Omega = 0^\circ\).

            To find \(\omega\) and \(i\), let us look at the following spherical triangle:

            \begin{figure}[H]
                \centering
                \includegraphics[angle=90, width=0.8\textwidth]{imagens/trigesfericoorbita.jpg}
                \caption{Diagram}
            \end{figure}

            From the polar form of the ellipse, \(r=\frac{a(1-e^2)}{1+e\cos\theta}\), where \(\theta\) is the true anomaly of the orbit, we can find \(\omega\). Notice that at the initial instant, \(t=0\), the probe is exactly at the point of Libra, so \(\theta = 2\pi-\omega\) as in the diagram.

            Thus, we have:

            \[d(0) = \frac{a(1-e^2)}{1+e\cos(2\pi-\omega)}\]

            Solving for \(\omega\):

            \[\omega = 2\pi-\cos^{-1}\left(\frac{a(1-e^2)}{d(0)e}-\frac{1}{e}\right) \approx 267.6^\circ \approx 270^\circ\]

            substituting the values found.

            To find the inclination, we will use the law of cosines:

            \[\cos(\theta-\omega)=\cos\lambda\cos b + \sin\lambda\sin b\cos(\pi/2) = \cos\lambda\cos b\]
            
            Note that \(\cos(\theta-\omega) \approx \cos(\theta-270^\circ) = \sin \theta\).

            Rewriting,

            \[\cos b = \frac{\sin\theta}{\cos\lambda}\]

            Using the relation obtained in the previous idea,

            \[\cos b  = \frac{r_T^2 + r^2 - d^2}{2 r_T r\cos (\lambda - \omega_\oplus t)}\]

            Equating these expressions and solving for \(\lambda\),

            \[\lambda = \tan^{-1} \left( 
                \frac{r_T^2 + r^2 - d^2}{2 r_T r \sin \theta \sin (\omega_\oplus t)} - \cot (\omega_\oplus t)
                \right)\]

            where:

            \[
            \theta = \arccos \left( \frac{a (1 - e^2)}{e r} - \frac{1}{e} \right)
            \] 
            
            With this, we can make the following table, with the values of \(b\) and \(\lambda\).

            \[
            \begin{array}{|c|c|}
            \hline
            \lambda \, (\text{Degrees}) & b \, (\text{Degrees}) \\ \hline
            0.000     & 0.00     \\ 
            17.495    & 9.847    \\ 
            45.905    & 22.521   \\ 
            78.492    & 29.499   \\ 
            101.508   & 29.499   \\ 
            134.095   & 22.521   \\ 
            163.504   & 9.847    \\ 
            180.000   & 0.000    \\ 
            197.495   & -9.847   \\ 
            216.005   & -18.747  \\ 
            270       & -30      \\ 
            333.435   & -14.478  \\ 
            17.495    & 9.847    \\ 
            78.492    & 29.499   \\ 
            143.994   & 18.747   \\ \hline
            \end{array}
            \]

            Returning to the spherical triangle and using the four-parts formula, we have:

            \[\cos\lambda\cos90^\circ = \sin\lambda\cot b -\sin90^\circ\cot i\]

            Simplifying:

            \[\tan b = \tan i \sin \lambda\]

            Note that this resembles a function of the type \(y= A + Bx\). Performing a linear regression on a calculator, we find \(B\approx 0.578\). Since \(\tan i = B\), we arrive at:

            \[\boxed{i \approx 30^\circ}\]

        \item Since the Sun is at the origin, the values of \(x'\) and \(y'\) can be calculated by:
            \[
            x' = r \sin \theta
            \]
            and
            \[
            y' = -r \cos \theta
            \]

            Writing five different values of \(x'\) and \(y'\) in a table, we obtain:

            \[
            \text{Table 7: Values of \(x'\) and \(y'\)}
            \]

            \[
            \begin{array}{|c|c|}
            \hline
            x'\text{ AU} & y'\text{ AU} \\ \hline
            0.167 & 0.000 \\ \hline
            0.084 & 0.478 \\ \hline
            -0.219 & 0.261 \\ \hline
            -0.203 & 0.074 \\ \hline
            0.000 & -0.100 \\ \hline
            \end{array}
            \]

            Notice that any values of \(x'\) and \(y'\), provided they are correct, may be used. Furthermore, it is important to increase the spacing between the values of \(x'\) and \(y'\) in order to reduce the errors associated with sketching the orbit. Based on this, we can plot the graph:

            \begin{figure}[H]
                \centering
                \includegraphics[width=0.7\linewidth]{imagens/graforbt.png}  
                \caption{Graph of the orbit} 
            \end{figure}

            Observing the graph, we see that the coordinates \((x', y', z')\) of the center of the ellipse are \(\boxed{(0;\ 0.2; \ 0)}\).

        \item Here, we will use a rotation matrix about the \(x\)-axis, given by:
        
            \[
            \begin{pmatrix}
            x \\
            y \\
            z
            \end{pmatrix}
            =
            \begin{pmatrix}
            1 & 0 & 0 \\
            0 & \cos \alpha & \sin \alpha \\
            0 & -\sin \alpha & \cos \alpha
            \end{pmatrix}
            \begin{pmatrix}
            x' \\
            y' \\
            z'
            \end{pmatrix}
            \] 

            For \(\alpha=i=30^\circ\):

            \[
            \begin{pmatrix}
            x \\
            y \\
            z
            \end{pmatrix}
            =
            \begin{pmatrix}
            1 & 0 & 0 \\
            0 & \cos 30^\circ & \sin 30^\circ \\
            0 & -\sin 30^\circ & \cos 30^\circ
            \end{pmatrix}
            \begin{pmatrix}
            x' \\
            y' \\
            z'
            \end{pmatrix}
            \]

            Then:
            \[
            x = x'
            \]
            \[
            y = y' \cos 30^\circ + z' \sin 30^\circ
            \]
            \[
            z = -y' \sin 30^\circ + z' \cos 30^\circ
            \]

            Substituting the numerical values, we obtain that the coordinates \((x, y, z)\) of the center of CR4b-2023 are given by:
            \[
            (0.000, 0.173, -0.100).
            \]

            Finally, we obtain that the distance between the center of Earth's orbit and the center of the probe's orbit is given by:
            \[
            d = \sqrt{\lvert x^2 + y^2 + z^2 \rvert}
            \]

            that is:
            \[
            \boxed{d = 0.200 \, \text{UA}}
            \]
        \end{alternativas}
    \end{pssolution*}
\end{pproblem}

\pts{5}
\begin{pproblem} (Iran Problem Set) 
    A particle of mass \(m\) is orbiting a massive object of mass \(M\). Show that the impulse required to rotate the orbit by an angle \(\eta\) about one of the foci is given by:
    \[
    \Delta v = \frac{2\mu e}{h}\sin\left(\frac{\eta}{2}\right),
    \]
    where:
    \begin{itemize}
        \item \(h\) is the angular momentum per unit mass;
        \item \(\mu = GM\), where \(G\) is the gravitational constant.
    \end{itemize}
    \begin{pssolution*}{}{}
        Drawing the orbit and the rotated orbit:

        \begin{figure}[H]
            \centering
            \includegraphics[width=0.7\linewidth]{imagens/esquemarotorb.png}
            \caption{Diagram of the rotated orbit.}
        \end{figure}

        Notice that the impulse can be applied only at the intersection points of the two orbits and must be radial, since the angular momentum must remain constant in order to keep the orbit identical. Let us define \(X\) as a point on the orbit before the impulse and \(X'\) as a point on the orbit after the rotation.

        Let us set the following definitions: \(F\) is the focus, \(P\) is the perihelion of the orbit, and \(A\) is the intersection point at which the impulse is applied.

        By symmetry, note that the angles \(AFP\) and \(AFF'\) are equal and therefore have the value \(\eta/2\).

        Since the two orbits have the same angular momentum, the tangential velocity at the point \(A\) must be the same for both orbits. Furthermore, we can see that the point \(A\) has true anomaly \(\theta = \eta/2\), which guarantees that the speed is also the same for both orbits.

        \begin{figure}[H]
            \centering
            \includegraphics[width=0.7\linewidth]{imagens/rotorbt1.png}
            \caption{Velocity vectors.}
        \end{figure}

        The only way for this to be possible is if the radial impulse is given by \(\Delta v = 2v_r\), where \(v_r\) is the radial velocity immediately before the impulse.

        Of course \(v_r^2 = v^2 - v_t^2\), where \(v_t\) is the tangential velocity. Starting from basic definitions:

        \[
        h^2 = \mu a (1 - e^2), \quad r = \frac{a(1 - e^2)}{1 - e\cos\theta}, \quad v_t = \frac{h}{r}, \quad v^2 = \mu \left(\frac{2}{r} - \frac{1}{a}\right).
        \]

        Working with these equations, we arrive at:

        \[
        v^2 = \frac{\mu^2}{h^2}(1 + 2e\cos\theta + e^2), \quad v_t^2 = \frac{\mu^2}{h^2}(1 + e\cos\theta)^2.
        \]

        Finally, substituting into \(v_r^2 = v^2 - v_t^2\), we obtain:

        \[
        v_r^2 = \frac{\mu^2}{h^2}(e^2 - e^2\cos^2\theta) \equiv \frac{\mu^2e^2}{h^2}\sin^2\theta.
        \]

        Substituting \(\Delta v = 2v_r\) and \(\theta = \eta/2\), we arrive at the desired result:

        \[
        \boxed{\Delta v = \frac{2\mu e}{h}\sin\left(\frac{\eta}{2}\right) \blacksquare}
        \]
    \end{pssolution*}
\end{pproblem}


\pts{3}
\begin{pproblem}
    In this question, we will study a simplified model of the effect that confirmed Einstein's theory of general relativity: gravitational lensing. The presence of a massive body curves spacetime, so that stars can serve as lenses in space. Some telescopes, such as the JWST, use this effect to photograph very distant galaxy clusters. One of the photographs taken by the JWST of a gravitational lens is shown below:
    \begin{figure}[H]
        \centering
        \includegraphics[width=0.7\linewidth]{imagens/fotoJWSTeinstenring.png}
        \caption{Source: National Geographic}
    \end{figure} 

    A simplified diagram of gravitational lensing can be seen below. This ``circle'' is known as the Einstein ring.
    \begin{figure}[H]
        \centering
        \includegraphics[width=0.9\linewidth]{imagens/anelE.png}
        \caption{Diagram of the gravitational lens}
    \end{figure}
    In the diagram, the point \(O\) represents the observer, \(L\) the massive body that acts as a lens, \(S\) the true position of the observed object, and \(S'\) its apparent position. \(\alpha\) is the impact parameter, \(\beta\) the angular separation between the observed body and the lens, \(\theta_E\) the angular radius of the Einstein ring, and \(\varphi\) the angular deflection caused by a massive body. The distance \(\overline{OL}\) equals \(D_L\) and the distance \(\overline{OP}\) equals \(D_S\), so that \(D_S-D_L = D_{LS}\).

    All the angles are very small, so that \(\sin x \approx \tan x \approx x\). Furthermore, because they lie at infinity, the lines \(\overline{OP}\), \(\overline{OS}\), and \(\overline{OS'}\) may be approximated as parallel.

    It is a known result of the theory of general relativity that the deflection of light caused by a massive body is given by:
    \[\varphi = \frac{4GM}{\alpha c^2}\]
    \begin{alternativas}
        \item In the limit \(\lim_{\alpha \rightarrow 0}\), find an expression for \(\theta_E\) in terms of the other distances and angles provided.
    
        \item The JWST telescope obtains images in the infrared, of wavelength \(\lambda_{IV}\). How large must its diameter be so that it can resolve an Einstein ring?
    \end{alternativas}

    \begin{pssolution*}{}{}
        \begin{alternativas}
            \item Note that, when \(\alpha \rightarrow 0\), \(\overline{PL}=\overline{AS}=D_{LS}\). Using \(\tan\theta = \sin\theta=\theta\), we have:
            
            \[(1)\ \overline{PS} = D_S\beta = \alpha, \ \ (2)\ \overline{SS'} = D_{SL}\varphi, \ \ (3)\ \overline{PS'} = D_S\theta_E,\] 
            \[(4)\ \overline{OL} = D_L = \overline{LA}/\theta_E = \alpha\theta_E, \ \ (5)\ \overline{LA}=\overline{PS}, \ \ (6)\ \overline{PS'} = \overline{PS}+\overline{SS'}\]

            Substituting \((1), \ (2), \ (3)\) into \((6)\), 

            \[D_S\theta_E=D_S\beta+(D_S-D_L)\varphi\]

            \[\theta_E-\beta = \frac{D_S-D_L}{D_S}\varphi\]

            Substituting the value of \(\varphi\):

            \[\theta_E-\beta = \frac{D_S-D_L}{D_S}\frac{4GM}{\alpha c^2}\]

            Using \((4)\), we have \(\alpha=D_L\theta_E\), and substituting into our equation:

            \[\theta_E-\beta = \frac{D_S-D_L}{D_S}\frac{4GM}{D_L\theta_E c^2}\]

            \[\theta_E^2-\beta\theta_E = \frac{D_S-D_L}{D_S}\frac{4GM}{c^2}\]

            The limit \(\alpha\rightarrow 0\) implies that \(\beta\rightarrow 0\), so:

            \[\boxed{\theta_E=\sqrt{\frac{D_S-D_L}{D_S}\frac{4GM}{c^2}}}\]

            \item Using the Rayleigh criterion \(\theta = R\frac{\lambda}{D}\), where \(R\approx 1.22\). Isolating \(D\) and substituting \(\theta = \theta_E\):
            
            \[\boxed{D =\sqrt{\frac{c^2R^2\lambda_{IV}^2}{4GM}\frac{D_S}{D_S-D_L}}}\]
        \end{alternativas}
    \end{pssolution*}
\end{pproblem}

\pts{4}
\begin{pproblem}
    In this exercise, we will study properties of the eccentricity of orbits and understand how it relates to energy and angular momentum. Denoting \(\mu = GM\) and \(h=L/m\), do what is asked in the following items.
    \begin{center}
        \textbf{Part I: The eccentricity vector}
    \end{center}
    \begin{alternativas}
        \item Write an expression for Newton's second law in vector form for a body of mass \(m\) orbiting a body of mass \(M\) at a distance \(r\).
        
        \item Take the cross product of your expression with \(\mathbf{h}\) and prove that:
        
        \[\frac{d}{dt}(\dot{\mathbf{r}}\times \mathbf{h}) = \frac{\mu}{r}\mathbf{v} - \frac{\mu \dot{r}}{r^2}\mathbf{r}\]

        where \(\dot{\mathbf{r}}\) represents the radial velocity and \(\mathbf{v}\) the velocity.

        \item We can write the right-hand side of the previous equation as:
        \[\frac{d}{dt}\left(A\mathbf{r}\right)\]

        Find the value of \(A\).

        \item Integrate both sides of the equality you found and take the scalar product of both sides with \(\mathbf{r}\). After that, solve for \(r\). Your result should be very similar to what is predicted by geometry alone:
        \[r = \frac{p}{1+e\cos\theta}\]

        Find the values of \(p\) and \(e\). \textbf{Do not forget the integration constant. Your answers may depend on it.}

        \item In the previous item, \(e\) is the eccentricity of the orbit. Returning to the equation initially obtained in item \(c\), after the integration, find a way to express \(\mathbf{e}\) as a vector of magnitude \(e\) that points directly toward the orbited body.
        
        \item Prove that the relation found in the previous item can be written as:
        \[\mu \mathbf{e} = (v^2-\frac{\mu}{r})\mathbf{r} - (\mathbf{r}\cdot\mathbf{v})\mathbf{v}\]

        Where does this vector point?
    \end{alternativas}
    \begin{center}
        \textbf{Part II: Determining orbital elements from the eccentricity vector}     
    \end{center}
    Let us define our coordinate system as follows. Consider a right-handed system, oriented so that the \(xy\) plane coincides with the equatorial plane and \(\mathbf{z}\) points toward the north pole. Define the nodal vector as \(\mathbf{n} = \mathbf{z} \times \mathbf{h}\).
    \begin{alternativas}
        \item Find an expression for \(\mathbf{h}\) in terms of the components of the vector \(\mathbf{r}\) and of the vector \(\mathbf{v}\). After that, find the vector \(\mathbf{n}\) in terms of the components of \(\mathbf{h}\).

        \item Draw a celestial sphere and represent on it: the equatorial plane and an arbitrary orbit (with inclination different from 0). On it, mark the following angles: \(i\), the orbital inclination, \(\Omega\) the longitude of the ascending node, \(\omega\) the argument of periastron, and \(\theta\), the true anomaly. 

        \item Find expressions for all of these angles in terms of \(\mathbf{e}\), \(\mathbf{h}\), \(\mathbf{n}\), \(\mathbf{r}\), and any of the vectors defined in our \(xyz\) coordinate system.
    \end{alternativas}

    \begin{pssolution*}{}{}
        \begin{center}
            \textbf{Part I}
        \end{center}
        \begin{alternativas}
            \item Using the vector form of the gravitational force, we have:
            \[-\frac{GMm}{r^3}\mathbf{r} = m\ddot{\mathbf{r}} \rightarrow \boxed{\ddot{\mathbf{r}}= -\frac{\mu}{r^3}\mathbf{r}}\]

            \item Taking the cross product of this expression with \(\mathbf{h}\), we have:
            
            \[\ddot{\mathbf{r}}\times \mathbf{h} = -\frac{\mu}{r^3}(\mathbf{r}\times\mathbf{h}) = \frac{\mu}{r^3}(\mathbf{h}\times\mathbf{r})\]

            Since \(\mathbf{h}\) is constant, the left-hand side of the equation can be written as \(d(\dot{\mathbf{r}}\times \mathbf{h})/dt\). For the right-hand side, let us write \(\mathbf{h} = \mathbf{r}\times \mathbf{v}\) and use the BAC-CAB rule, which states that \((A \times B) \times C = B(A\cdot C) - C(A \cdot B)\). Thus:

            \[\frac{d}{dt}(\dot{\mathbf{r}}\times \mathbf{h}) = \frac{\mu}{r^3}((\mathbf{r}\times\mathbf{v})\times \mathbf{r}) = \frac{\mu}{r^3}(r^2\mathbf{v} - \mathbf{r}(\mathbf{r}\cdot \mathbf{v}))\]

            Note that \(\mathbf{r}\cdot\mathbf{v}\) is the product of the components of \(\mathbf{r}\) and \(\mathbf{v}\) that lie in the same direction. Thus, \(\mathbf{r}\cdot\mathbf{v} = r\dot{r}\). Substituting into the previous expression:

            \[\boxed{\frac{d}{dt}(\dot{\mathbf{r}}\times \mathbf{h}) = \frac{\mu}{r}\mathbf{v} - \frac{\mu \dot{r}}{r^2}\mathbf{r}}\]

            \item Expanding the expression given in the statement:
            
            \[\frac{d}{dt}(A\mathbf{r}) = \mathbf{r}\frac{d}{dt}A + A \mathbf{v}\]

            Comparing the expressions, it is easy to see that \(\boxed{A = \mu/r}\).

            \item Using the information obtained in the previous items, we have:
            \[\frac{d}{dt}(\dot{\mathbf{r}}\times \mathbf{h}) = \mu\frac{d}{dt}\left(\frac{\mathbf{r}}{r}\right)\]

            Multiplying both sides by \(dt\) and integrating, we have:

            \[\dot{\mathbf{r}}\times\mathbf{h} = \frac{\mu}{r}\mathbf{r} + \mathbf{B}\]

            where \(B\) is a constant vector of integration. Now, taking the scalar product with \(\mathbf{r}\):

            \[\mathbf{r}\cdot(\dot{\mathbf{r}}\times \mathbf{h}) = \mu r + \mathbf{r}\cdot\mathbf{B}\]

            Using \(a\cdot(b\times c) = (a\times b)\cdot c\) and defining \(\nu\) as the angle between \(\mathbf{r}\) and \(\mathbf{B}\), we have:

            \[h^2 = \mu r + rB\cos\nu\]

            Solving for \(r\), we arrive at:

            \[r = \frac{h^2/\mu}{1+(B/\mu)\cos\nu}\]

            Thus:

            \[\boxed{p = \frac{h^2}{\mu}, \ \ \ \ e = \frac{B}{\mu}}\]

            \item Using \(\mathbf{e} = \mu \mathbf{B}\) and returning to item \(d)\), we have:
            \[\dot{\mathbf{r}}\times \mathbf{h} = \frac{\mu}{r}\mathbf{r}+\mu \mathbf{e}\]

            Solving for \(\mathbf{e}\) and noting that \(\dot{\mathbf{r}}\times\mathbf{h} = \mathbf{v}\times\mathbf{h}\), we have:

            \[\boxed{\mathbf{e} = \frac{\mathbf{v}\times\mathbf{h}}{\mu} - \frac{\mathbf{r}}{r}}\]

            \item Writing \(\mathbf{h} = \mathbf{r}\times\mathbf{v}\), we have:
            \[\mathbf{e} = \frac{\mathbf{v}\times(\mathbf{r}\times\mathbf{v})}{\mu} - \frac{\mathbf{r}}{r} \]

            Multiplying both sides of the expression by \(\mu\), and using BAC-CAB, we have:

            \[\mu\mathbf{e} = v^2\mathbf{r} - \mathbf{v}(\mathbf{v}\cdot\mathbf{r}) - \frac{\mu}{r}\mathbf{r}\]

            Rearranging the terms:

            \[\boxed{\mu \mathbf{e} = (v^2-\frac{\mu}{r})\mathbf{r} - (\mathbf{r}\cdot\mathbf{v})\mathbf{v}}\]

            Thus, this vector points toward periastron.
        \end{alternativas}
        \begin{center}
            \textbf{Part II: Determining orbital elements from the eccentricity vector}     
        \end{center}
        \begin{alternativas}
            \item The vector $\mathbf{h}$ is defined by $\mathbf{h} = \mathbf{r}\times\mathbf{v}$. By definition,
            
            $$\mathbf{h}=\mathbf{r}\times\mathbf{v} = 
            \begin{vmatrix}
                \hat{x} & \hat{y} & \hat{z} \\
                r_x & r_y & r_z \\
                v_x & v_y & v_z \\
            \end{vmatrix}
            = (r_yv_z - r_zv_y)\hat{x} + (r_zv_x-r_xv_z)\hat{y} + (r_xv_y-r_yv_z)\hat{z}   $$

            Similarly, $\mathbf{n} = \mathbf{z}\times\mathbf{h}$
            $$\mathbf{n} = 
            \begin{vmatrix}
                \hat{x} & \hat{y} & \hat{z} \\
                0 & 0 & 1 \\
                h_x & h_y & h_z \\
            \end{vmatrix}
             = -h_y\hat{x} + h_x\hat{y}$$

             \item The drawing should resemble the following figure:
             \begin{figure}[H]
                \centering
                \includegraphics[width=0.7\linewidth]{imagens/elementosorbitais.png}
                \caption{Source: Wikipedia}
            \end{figure}

            \item To orient ourselves, let us create a coordinate system in which the north pole corresponds to the $z$-axis and the $x$-axis points in the direction of the vernal point. Starting with the inclination, it is the angle between the north pole and the plane perpendicular to the orbit. Thus
            $$\mathbf{z}\cdot\mathbf{h} = \cos i |z||h|\rightarrow i = \cos^{-1}\left(\frac{\mathbf{z}\cdot\mathbf{h}}{|z||h|}\right) $$

            Using similar arguments, we can obtain

            $$\Omega = \cos^{-1}\left(\frac{\mathbf{x}\cdot\mathbf{n}}{|x||n|}\right)$$
            $$\omega = \cos^{-1}\left(\frac{\mathbf{n}\cdot\mathbf{e}}{|n||e|}\right)$$
            $$\theta = \cos^{-1}\left(\frac{\mathbf{e}\cdot\mathbf{r}}{|e||r|}\right)$$
            
        \end{alternativas}
    \end{pssolution*}
\end{pproblem}

\pts{5}
\begin{pproblem}\(\star\star\star\) (Adapted from IPhO 2018) One of the most interesting effects of general relativity is the emission of gravitational waves (GWs) by close binaries. One consequence of this is the loss of energy of the system. In this question, we will study this effect in depth.

\begin{alternativas}
    \item Consider a system formed by stars 1 and 2, with masses \(M_i\) and at distances \(r_i\) from the center of mass. Using Newton's second law, one can show that:
    \[\mathbf{a_1} = -\alpha \frac{\mathbf{r_1}}{r_1^n}\]
    where \(\mathbf{a_1}\) is the acceleration vector of star 1. Find the value of \(\alpha = \alpha (G, M_1, M_2)\) (read as \(\alpha\) in terms of \(G, M_1, M_2\)) and of \(n\).

    \item The total energy of the binary system can be expressed as:
    \[E = A(\omega, \mu, L) - \frac{GM\mu}{L}\]

    where \(\mu\) is the reduced mass of the system, \(M = M_1+M_2\), and \(L = r_1+r_2\). Find the value of \(A\).

    \item The previous expression can be simplified to:
    \[E = \frac{\beta GM\mu}{L}\]

    where \(\beta\) is a dimensionless constant. Find its value.

    The correct theory of gravity, general relativity, was formulated by Einstein in 1915 and predicts that gravity propagates at the speed of light. The messengers that carry information about the interaction are called gravitational waves (GWs). GWs are emitted whenever masses are accelerated, causing the system of masses to lose energy.

    Consider a system of two point particles, isolated from the rest of the Universe. Einstein proved that, for sufficiently small speeds, the emitted waves: 1) have a frequency twice the orbital frequency; 2) can be characterized by a luminosity, that is, an emitted power \(\mathcal{P}\), which is dominated by the formula:

    \[
    \mathcal{P} = \frac{G}{5c^5} \sum_{i=1}^3 \sum_{j=1}^3 \left( \frac{d^3 Q_{ij}}{dt^3} \right) \left( \frac{d^3 Q_{ij}}{dt^3} \right).
    \]

    Here, \( c \) is the speed of light \( c \simeq 3 \times 10^8 \, \mathrm{m/s} \). For a system of two point particles orbiting in the \( x-y \) plane, \( Q_{ij} \) is the following table (\( i, j \) indicate the row/column number):

    The components of \( Q_{ij} \) are given by:

    \[
    Q_{11} = \sum_{A=1}^2 \frac{M_A}{3} \left( 2x_A^2 - y_A^2 \right),
    \]

    \[
    Q_{22} = \sum_{A=1}^2 \frac{M_A}{3} \left( 2y_A^2 - x_A^2 \right),
    \]

    \[
    Q_{33} = -\sum_{A=1}^2 \frac{M_A}{3} \left( x_A^2 + y_A^2 \right),
    \]

    \[
    Q_{12} = Q_{21} = \sum_{A=1}^2 M_A x_A y_A.
    \]

    And \(Q_{i,j} = 0\) for all other possibilities. Here, \(x_A, y_A\) are the positions of star \(A\), with \(A=1,\ 2\), in a Cartesian plane with the center of mass at the origin.

    \item Write the coordinates (\(x_1, y_1\)) and (\(x_2, y_2\)) in terms of \(r_1, r_2, \omega\), and \(t\), with \(t\) the time elapsed since some specific point of the orbit.
    
    \item To calculate the dissipated power, the formula contains a quadruple sum. Solving this expression by brute force is quite complicated and time-consuming. Fortunately, there is a nicer and more elegant way to solve it. To do so, begin by writing \(Q_{i,j}\) as a 3x3 matrix, of the following form:
    
    \[Q = A\left(\begin{matrix}
        a_{11} & a_{12} & a_{13} \\
        a_{21} & a_{22} & a_{23} \\
        a_{31} & a_{23} & a_{33} \\
    \end{matrix}\right)\]

    Find the coefficient \(A = A(\mu, L)\) and complete the matrix \(Q\).

    \item From the previous item, you should obtain:
    \[Q_{ii} = A(b_i + j_i\cos(k t)), \ \ \ Q_{ij}^{i\neq j}= A(p_{ij} \sin(kt)) \]

    Find the values of \(b_i\), \(j_i\), \(p_{ij}\), and \(k\).

    \item You are now able to solve for \(\mathcal{P}\) more smoothly. The expression can be simplified to:
    \[\mathcal{P} = \xi \frac{G}{c^5}\mu^2L^4\omega^6\]
    Find the numerical value of \(\xi\).

    \textbf{HINT: } The double sum:

    \[\sum_{i=1}^N\sum_{j=1}^N (A_{ij})(A_{ij})\]

    where \(A\) is an \(N\times N\) matrix, can be simplified to:

    \[\sum_{i=1}^N\sum_{j=1}^N A_{ij}^2\]

    which represents the sum of the squares of all the elements of the matrix \(A\).

    \item If there were no emission of GWs, the system would remain in equilibrium indefinitely, but because of that emission the energy of the system is not conserved, so that there is a change in the angular velocity of the system, \(\omega\). The formula for the time variation of \(\omega\) has the following form:
    \[\left(\frac{d\omega}{dt}\right)^3 = (3\xi )^3 \frac{\omega^{11}}{c^{15}}(GM_C)^5\]
    
    where \(M_c\) is the so-called \textit{chirp} mass and \(M_c = M_c(M, \mu)\). Find a formula for \(M_c\).

    \item Using the information obtained above, relate the orbital angular velocity \(\omega \) to the gravitational-wave frequency \( f_{\text{GW}} \). Given that, for any smooth function \( F(t) \) and \( a \neq 1 \):

    \[
    \frac{dF(t)}{dt} = \chi F(t)^a \implies F(t)^{1-a} = \chi (1-a)(t_0 - t),
    \]
    
    where \( \chi \) is a constant and \( t_0 \) is an integration constant, show that the equation of the previous item implies that the gravitational-wave frequency is:
    
    \[
    f_{\text{GW}}^{-8/3} = 8\pi^{8/3} \xi \left( \frac{GM_c}{c^3} \right)^{(2/3)+p} (t_0 - t)^{2-p},
    \]
    
    and determine the constant \( p \).

    On September 14, 2015, the event GW150914 was recorded by the LIGO detectors, consisting of two L-shaped arms, each 4 km long. These arms changed their relative length according to the figure below. The detector arms respond linearly to a passing gravitational wave, and the response pattern mimics the wave. This wave was created by two black holes on nearly circular orbits; the loss of energy by gravitational radiation caused the orbit to shrink and, eventually, the black holes to collide. The collision point corresponds, approximately, to the peak of the signal after point D in the figure below.

    \begin{figure}[H]
        \centering
        \includegraphics[width=0.7\linewidth]{imagens/grafico OG 1.png}
        \caption{Strain, that is, the relative change in the length of each arm, in the LIGO H1 detector. The horizontal axis represents time, and the points A, B, C, and D correspond to \(t = 0.000\), \(0.009\), \(0.034\), and \(0.040\) seconds, respectively.}
    \end{figure}

    \item From the figure, estimate \( f_{GW}(t) \) at:

    \[
    t_{AB} = \frac{t_B + t_A}{2} \quad \text{and} \quad t_{CD} = \frac{t_D + t_C}{2}.
    \]
    
    Assuming that the equation for \(f_{GW}\) is valid up to the collision (which, strictly speaking, is not true) and that the two objects have equal masses, estimate the chirp mass, \( M_c \), and the total mass of the system, in terms of solar masses \( M_\odot \simeq 2 \times 10^{30} \, \text{kg}\).

    \item Estimate the minimum orbital separation between the two objects at \( t_{CD} \). 
    Then estimate a maximum size for each object, \( R_{\text{max}} \). 
    Obtain \( \frac{R_\odot}{R_{\text{max}}} \) in order to compare this size with the radius of our Sun, \( R_\odot \simeq 7 \times 10^5 \, \text{km} \). 
    Also estimate their linear orbital speed at the same instant, \( v_{\text{col}} \), comparing it with the speed of light, \( \frac{v_{\text{col}}}{c} \).
\end{alternativas}

\begin{pssolution*}{}{}
    \begin{alternativas}
        \item From Newton's second law,
        \[-\frac{GM_1M_2}{(r_1+r_2)^3}\mathbf{d} = M_1\mathbf{a_1}\]

        where \(\mathbf{d}\) is a vector that leaves star two and goes to star one. Therefore, \(|\mathbf{d}| = r_1+r_2\). Using the center-of-mass theorem, we have

        \[r_2 = \frac{M_1r_1}{M_2}\]

        Substituting and simplifying,

        \[\mathbf{a_1} = -\frac{GM_2}{r_1^3(1+\frac{M_1}{M_2})^3}\mathbf{d}\]

        Now working with the vector \(\mathbf{d}\), note that its direction is the same as that of \(\mathbf{r_1}\), since the line joining 1 and 2 necessarily passes through the CM of the system. Thus,

        \[\mathbf{d} = |\mathbf{d}|\hat{\mathbf{r_1}} = (r_1+r_2)\hat{\mathbf{r_1}} = (1+M_1/M_2)\mathbf{r_1}\]

        Substituting,

        \[\mathbf{a_1} = -\frac{GM_2}{(1+M_1/M_2)^2}\frac{\mathbf{r_1}}{r_1^3}\]

        With this, we can conclude that

        \[\boxed{\alpha = \frac{GM_2}{(1+M_1/M_2)^2}, \ \ \ n = 3 \ }\]

        \item The energy of the system is given by
        \[E = \frac{M_1v_1^2+M_2v_2^2}{2}-\frac{GM_1M_2}{L}\]

        Focusing on the first term and using \(v_A = r_A\omega\), we have

        \[M_1v_1^2+M_2v_2^2 = (M_1r_1^2 + M_2r_2^2)\omega^2 \]

        By the center-of-mass theorem

        \[(M_1r_1-M_2r_2)^2=0 \rightarrow M_1r_1^2 +M_2r^2_2 = \frac{M_1M_2}{M_1+M_2}L^2\equiv \mu L^2\]

        For the second term, it suffices to multiply the numerator and the denominator by \((M_1+M_2)\). Thus,

        \[E = \frac{\mu L^2}{2} - \frac{G}{L}\frac{M_1M_2}{M_1+M_2}(M_1+M_2) \equiv \frac{\mu L^2\omega^2}{2} + \frac{GM\mu }{L}\]

        Thus, \(\boxed{A = \frac{\mu L^2\omega^2}{2}}\).

        \item By Kepler's third law,
        
        \[\omega^2 = \frac{GM}{L^3}\]

        Substituting,

        \[E = \frac{GM\mu}{2L} - \frac{GM\mu}{L} = -\frac{GM\mu}{L}\]

        thus concluding that \(\boxed{\beta = -1/2}\).

        \item Let \(\theta\) be the angle that the vector \(r_1\) makes with the \(x\)-axis. Thus,
        \[(x_1, \ y_1) = r_1(\cos\theta, \ \sin\theta)\]

        The vector \(r_2\) must make an angle of \(\pi + \theta\). Using \(\sin(\pi+\theta) = -\sin(\theta)\) and \(\cos(\pi+\theta) = -\cos\theta\),

        \[(x_2, \ y_2) = -r_2(\cos\theta, \sin\theta)\]

        Defining \(t\) as the time elapsed since the stars crossed the \(x\)-axis, one has \(\theta = \omega t\). Therefore,

        \[\boxed{(x_1, \ y_1) = r_1 (\cos(\omega t), \ \sin(\omega t)) \ \ \ (x_2, \ y_2) = -r_2 (\cos(\omega t), \ \sin(\omega t))}\]

        \item The matrix will have the following form,
        
        \[Q = \left(\begin{matrix}
            Q_{11} & Q_{12} & Q_{13} \\
            Q_{21} & Q_{22} & Q_{23} \\
            Q_{31} & Q_{32} & Q_{33} \\
        \end{matrix}\right)\]

        Invoking the formulas from the statement,
        \[
        Q_{11} = \sum_{A=1}^2 \frac{M_A}{3} \left( 2x_A^2 - y_A^2 \right),
        \]

        \[
        Q_{22} = \sum_{A=1}^2 \frac{M_A}{3} \left( 2y_A^2 - x_A^2 \right),
        \]

        \[
        Q_{33} = -\sum_{A=1}^2 \frac{M_A}{3} \left( x_A^2 + y_A^2 \right),
        \]

        \[
        Q_{12} = Q_{21} = \sum_{A=1}^2 M_A x_A y_A.
        \]

         For \(Q_{11}\), we have

        \[Q_{11} = \frac{M_1}{3}(2x_1^2-y_1^2) + \frac{M_2}{3}(2x_2^2-y_2^2)\]

        Substituting \(x\) and \(y\),

        \[Q_{11} = \frac{M_1}{3}r_1^2(2\cos^2(\omega t) - \sin^2(\omega t)) + \frac{M_2r^2_2}{3}(2\cos^2(\omega t) - \sin^2(\omega t))\]

        \[Q_{11} = \frac{1}{3}(2\cos^2(\omega t) - \sin^2(\omega t))(M_1r_1^2 + M_2r_2^2)\]  
        \[Q_{11} = \frac{\mu L^2}{3}(2\cos^2(\omega t) - \sin^2(\omega t))\]

        For \(Q_{22}\),

        \[Q_{22} = \frac{M_1r_1^2}{3}(2\sin^2(\omega t) - \cos^2(\omega t)) + \frac{M_2r_2^2}{3}(2\sin^2(\omega t) - \cos^(\omega t))\]

        \[Q_{22} = \frac{\mu L^2}{3}(2\sin^2(\omega t) - \cos^2(\omega t))\]


        For \(Q_{33}\)

        \[Q_{33} = -\frac{M_1r_1^2}{3}(\cos^2(\omega t) + \sin^2(\omega t))-\frac{M_2r_2^2}{3}(\cos^2(\omega t) + \sin^2(\omega t))\]

        \[Q_{33} = -\frac{\mu L^2}{3}\]

        Finally, for \(Q_{12} = Q_{21}\),

        \[Q_{12} = Q_{21} = M_1r_1^2\cos(\omega t) \sin(\omega t) + M_2r_2^2\cos(\omega t)\sin(\omega t)\]

        \[Q_{12} = Q_{21} = \mu L^2 \sin(\omega t)\cos(\omega t) = \frac{\mu L^2}{2}\sin(2\omega t)\]

        Since the remaining terms are \(0\), we can choose as \(A\) the value \(\mu L^2/2\), so our matrix \(Q\) takes the form

        \[Q = \frac{\mu L^2}{2}\left(\begin{matrix}
            \frac{4}{3}\cos^2(\omega t) - \frac{2}{3}\sin^2(\omega t)  & \sin(2\omega t) & 0 \\
            \sin(2\omega t) & \frac{4}{3}\sin^2(\omega t) - \frac{2}{3}\cos(\omega t) & 0 \\
            0 & 0 & -\frac{2}{3} \\
            \end{matrix}\right)\]

        For the terms \(a_{11}\) and \(a_{22}\), we will use trigonometric identities, starting with \(a_{11}\),
        \[\frac{4}{3}\cos^2(\omega t) - \frac{2}{3}\sin^2(\omega t) = \frac{2}{3}(2\cos^2(\omega t) - \sin^2(\omega t))\]

        Using \(\cos^2(\omega t) = 1-\sin^2(\omega t)\), we have

        \[ = \frac{2}{3}\left(2-3\sin^{2}(\omega t)\right)\]

        Using the double-angle form, we have

        \[\cos(2x) = \cos^2x -\sin^2x = 1-2\sin^2x \rightarrow \sin^2x = \frac{1-\cos(2x)}{2}\]

        Making this substitution,

        \[=\frac{2}{3}\left(2 - \frac{3}{2} + \frac{3}{2}\cos(2\omega t)\right) = \frac{1}{3}+\cos(2\omega t)\]

        Using the same process for \(a_{22}\), we find
        \[a_{22} = \frac{1}{3}-\cos(2\omega t)\]

        so that our matrix can be simplified to


        \[\boxed{Q = \left(\begin{matrix}
            \frac{1}{3}+\cos(2\omega t) & \sin(2\omega t) & 0 \\
            \sin(2\omega t) & \frac{1}{3} - \cos(2\omega t) & 0 \\
            0 & 0 & -\frac{2}{3} \\
        \end{matrix}\right)}\]

        \item Analyzing the previous matrix, we have
        \[\boxed{A = \frac{\mu L^2}{2}}\]
        and
        \[\boxed{b_1 = b_2 = \frac{1}{3}, \ \ b_3 = -\frac{2}{3}, \ \ j_1 = 1, \ \ j_2 = -1, \ \ j_3 = 0, \ \ p_{12}=p_{21} = 2, \ \ k= 2}\]

        \item The formula provided for the power was
        
        \[
        \mathcal{P} = \frac{G}{5c^5} \sum_{i=1}^3 \sum_{j=1}^3 \left( \frac{d^3 Q_{ij}}{dt^3} \right) \left( \frac{d^3 Q_{ij}}{dt^3} \right).
        \]

        To begin, let us differentiate the matrix \(Q\) with respect to time,

        \[\frac{d^3 Q}{dt^3} = \frac{\mu L^2}{2}\frac{d^3}{dt^3}\left(\begin{matrix}
            \frac{1}{3}+\cos(2\omega t) & \sin(2\omega t) & 0 \\
            \sin(2\omega t) & \frac{1}{3} - \cos(2\omega t) & 0 \\
            0 & 0 & -\frac{2}{3} \\
        \end{matrix}\right)\]

        To differentiate a matrix, it suffices to differentiate all of its elements. Doing this, we arrive at

        \[\frac{d^3Q}{dt^3} = \frac{\mu L^2}{2}\left(\begin{matrix}
            8\omega^3 \sin(2\omega t) & -8\omega ^3\cos(2\omega t) & 0 \\
            -8\omega^3\cos(\omega t) & -8\omega^3\sin(\omega t) & 0 \\
            0 & 0 & 0 \\ 
        \end{matrix}\right)\]

        Taking that factor outside the matrix,

        \[\frac{d^3Q}{dt^3} = 4\mu\omega^3L^2\left(\begin{matrix}
            \sin(2\omega t) & - \cos(2\omega t) & 0 \\
            -\cos(2\omega t) & - \sin(\omega t) & 0 \\
            0 & 0 & 0 \\
        \end{matrix}\right)\]

        It was given as a hint that the term

        \[\sum_{i=1}^3 \sum_{j=1}^3 \left( \frac{d^3 Q_{ij}}{dt^3} \right) \left( \frac{d^3 Q_{ij}}{dt^3} \right)\]

        corresponds to the sum of the squares of the entries of the matrix, so

        \[\sum_{i=1}^3 \sum_{j=1}^3 \left( \frac{d^3 Q_{ij}}{dt^3} \right) \left( \frac{d^3 Q_{ij}}{dt^3} \right) = (4\mu \omega^3 L^2)^2(\sin^2(2\omega t) + 2 \cos^2(\omega t))\]

        Thus,

        \[\mathcal{P} = \frac{G}{5c^5} \sum_{i=1}^3 \sum_{j=1}^3 \left( \frac{d^3 Q_{ij}}{dt^3} \right) \left( \frac{d^3 Q_{ij}}{dt^3} \right) = \frac{G}{5c^5}(32\mu^2\omega^6 L^4)\]

        Rearranging,

        \[\mathcal{P} = \frac{32}{5}\frac{G\mu^2\omega^6L^4}{c^5}\]

        Comparing this expression with the one given in the statement, we have

        \[\boxed{\xi = \frac{32}{5}}\]

        \item \(\mathcal{P}\) represents the power being dissipated by the GWs, so
        \[\mathcal{P} = -\frac{dE}{dt}\]

        Equating \(E\) to the orbital energy,

        \[\mathcal{P} = \frac{GM\mu}{2} \frac{d}{dt}\frac{1}{L}\]

        We can write \(\omega\) as a function of \(L\). To do so, it suffices to use Kepler's third law,


        \[\omega^2 = \frac{GM}{L^3} \rightarrow \frac{1}{L }= (GM)^{-1/3}\omega^{2/3}\]

        Substituting,

        \[\mathcal{P} = \frac{(GM)^{2/3}\mu}{2}\frac{d}{dt}\omega^{2/3} = \frac{(GM)^{2/3}\mu}{2}\frac{d}{d\omega}\omega^{2/3}\frac{d\omega}{dt}\]

        \[= \frac{(GM)^{2/3}\mu}{3\omega^{1/3}}\frac{d\omega}{dt}\]

        Equating this to the expression obtained in the previous item,

        \[\frac{\xi G\mu^2\omega^6L^4}{c^5} = \frac{(GM)^{2/3}\mu}{3\omega^{1/3}}\frac{d\omega}{dt}\]

        Solving for \(d\omega/dt\),
        
        \[\frac{d\omega }{dt} = (3\xi)\frac{G\mu \omega^{19/3}L^4}{(GM)^{2/3}c^5}\]

        Substituting \(L = \left(\frac{GM}{\omega^2}\right)^{1/3}\), we have

        \[\frac{d\omega}{dt} = (3\xi)\frac{G\mu \omega^{19/3}}{(GM)^{2/3}c^5}\left(\frac{GM}{\omega^2}\right)^{4/3}\]

        Cubing,

        \[\left(\frac{d\omega}{dt}\right)^{3} = (3\xi)^3\frac{G^3\mu^3\omega^{19}}{G^2M^2c^{15}}\frac{G^4M^4}{\omega^8}\]

        Simplifying,

        \[\boxed{\left(\frac{d\omega}{dt}\right)^3 = (3\xi)^3\frac{\omega^{11}}{c^{{15}}}(G^5\mu^3M^2)}\]

        Analyzing the expression in the statement, we have

        \[\boxed{M_c = (\mu^3 M^2)^{1/5}}\]

        \item Using the fact given in the statement,
        
        \[\frac{d F(t)}{dt} = \chi F(t)^a \rightarrow F(t)^{1-a} = \chi (1-a)(t_0-t)\]

        we can find an expression for \(\omega(t)\).

        \[\frac{d\omega}{dt} = \frac{3\xi (GM_c)^{3/5}}{c^5}\omega^{11/3} \rightarrow \omega^{1-11/3} = \frac{3\xi (GM_c)^{3/5}}{c^5}(1-11/3)(t_0-t)\]

        Simplifying,

        \[\omega ^{-8/3} = -\frac{3\xi (GM_c)^{3/5}}{c^5}\frac{8}{3}(t_0-t)\]

        \[\omega^{-8/3} = \frac{8\xi (GM_c)^{3/5}(t-t_0)}{c^5}\]

        To find an expression for \(f_{GW}\), we must use

        \[f_{GW} = 2f_{orbt} = \frac{\omega}{\pi}\rightarrow \omega = \pi f_{GW}\]

        Making this substitution, we arrive at


        \[f_{GW}^{-8/3} = \frac{8\pi^{8/3}(GM_c)^{3/5}(t-t_0)}{c^5}\]

        Comparing with the form in the statement, we obtain

        \[\boxed{p = 1}\]

        \item From the figure, we take the two $\Delta t$ as half-periods. Thus, the cyclic gravitational-wave frequency is $f_{GW} = 1/(2 \Delta t)$. The four given points then allow us to calculate the frequency at the mean time of the two intervals as:

        \[
        \begin{array}{|c|c|}
        \hline
        t \, (\text{s}) & f_{GW} \, (\text{Hz}) \\
        \hline
        t_{AB} = 0.0045 & f_{GW}(t_{AB}) = (2 \times 0.009)^{-1} \\
        t_{CD} = 0.037 & f_{GW}(t_{CD}) = (2 \times 0.006)^{-1} \\
        \hline
        \end{array}
        \]
        
        Now we have two pairs of values $(f_{GW}, t)$ for two unknowns $(t_0, M_c)$. In terms of $t_{AB}$ and $t_{CD}$, and dividing the two equations, we obtain:
        
        \[
        t_0 = \frac{A t_{CD} - t_{AB}}{A - 1}, \quad A = \left( \frac{f_{GW}(t_{AB})}{f_{GW}(t_{CD})} \right)^{-8/3}.
        \]
        
        Substituting the numerical values, $A \simeq 2.95$ and $t_0 \simeq 0.054 \, \text{s}$. Now we can use the equation for \(f_{GW}\) at either $t_{AB}$ or $t_{CD}$ and determine $M_c$. One obtains the chirp mass:
        
        \[
        M_c \simeq 6 \times 10^{31} \, \text{kg} \simeq 30 \, M_\odot\]
        
        Thus, the total mass $M$ is:
        
        \[
        M = 4^{3/5} M_c \simeq 69 \, M_\odot
        \]
        
        This result is, in fact, remarkably close to the best estimates that use the full theory of general relativity! [Even though the real objects do not have exactly equal masses and the theory we have just used is not valid very close to the collision.]
    
        \item Kepler's law gives $L = (GM / \Omega^2)^{1/3}$. The second pair of points highlighted on the graph corresponds to the cycle before the merger. Thus, using \(f_{GW} = \omega/\pi\),

        \[
        \omega_{t_{CD}} \sim 2.6 \times 10^2 \, \text{rad/s}.
        \]
        
        Next, using the total mass obtained previously:
        
        \[
        L \sim 5 \times 10^2 \, \text{km}
        \]
        
        Thus, these objects have a maximum radius of $R_{\text{max}} \sim 250 \, \text{km}$. Therefore, they have more than 30 times as much mass, and:
        
        \[
        \frac{R_\odot}{R_{\text{max}}} \sim 3 \times 10^3,
        \]
        
        they are 3000 times smaller than the Sun!
        
        The linear speed is:
        
        \[
        v_{\text{col}} = \frac{L}{2} \omega \simeq 7 \times 10^4 \, \text{km/s}. 
        \]
        
        They are moving at more than 20\% of the speed of light!
        
        \[
        \boxed{
        \begin{aligned}
        \quad & L_{\text{col}} \sim 5 \times 10^2 \, \text{km}, \\
        & \frac{R_\odot}{R_{\text{max}}} \sim 3 \times 10^3, \\
        & \frac{v_{\text{col}}}{c} \sim 0.2.
        \end{aligned}
        }
        \]
    \end{alternativas}
    
\end{pssolution*}
\end{pproblem}

\pts{3}
\begin{pproblem}(Barra 2024) During his stay at Hotel Fazenda Ribeirão, Hugo was observing the star Plo I with his telescope. Comparing its spectrum with those of other stars, he discovers that Plo I is a neutron star and estimates its mass to be \( M_1 = 2M_{\odot} \). He also observes that Plo I has sinusoidal variations in its radial velocity over time and concludes that Plo I must be part of a binary system with another celestial body, Plo II. Observing the velocity curve of Plo I, Hugo also determines that Plo I has a circular orbit of period \( P \) and maximum radial velocity \( K \). However, he does not know the mass \( M_2 \) of Plo II, nor the inclination \( i \) of the binary's orbital plane with respect to the plane of the sky, shown in the following figure.

    \begin{figure}[H]
        \centering
        \includegraphics[width=0.8\linewidth]{imagens/figurabarra.png}
        \caption{Representation of the orbit}
    \end{figure}

    \begin{alternativas}
        \item From Kepler's third law and Newton's second law, find an expression for the \textbf{mass function} of the binary system formed by Plo I and Plo II in terms of \( P \), \( K \), and universal constants. You may use that the mass function is given by:
        \[
            f = \frac{M_2^3 \sin^3 i}{(M_1 + M_2)^2}
        \]
    
        \item Using his measurements, Hugo calculates that the mass function of the binary is \( f = 46.2M_{\odot} \). We wish to find the minimum value \( M_{2,\text{min}} \) of the mass of Plo II, knowing that it corresponds to the case \( i = 90^\circ \). Assume \( M_{2,\text{min}} > M_1 \). From this, determine \( M_{2,\text{min}} \). Is your result consistent with the hypothesis? Answer YES or NO, justifying with calculations.
        
        \item Based on your answer to the previous item, conclude: is it more likely that Plo II is an exoplanet, a white dwarf, a neutron star, or a black hole?
    \end{alternativas}
    \begin{pssolution*}{}{}
        \begin{alternativas}
            \item First, note that \(v_{r,\ \mathrm{apparent}} = |\mathbf{\omega}\times\mathbf{r}| = \omega r \sin i = v_r \sin i\), so \(K = v_r\sin i = \frac{2\pi a_1}{T}\sin i\). Defining \(a=a_1+a_2\) and \(M = M_1+M_2\), where \(a_i\) and \(M_i\) represent the distance between star \(i\) and the CM and the mass of star \(i\), we have Kepler's third law as
            \[\frac{a^3}{T^2} = \frac{GM}{4\pi^2}\]

            \[a^3 = \frac{GM}{4\pi^2}T^2 \rightarrow T = 2\pi \sqrt{\frac{a^3}{GM}}\]

            Substituting $T = \frac{2\pi a_1}{v_1}$

            $$(a_1+a_2)^3= \frac{G(M_1+M_2)a_1^2}{v_1^2}$$

            $$(1+a_2/a_1)^2 = \frac{G(M_1+M_2)}{(a_1+a_2)v_1^2}$$

            Also, using the center-of-mass theorem,

            $$\frac{a_1}{a_2} = \frac{M_2}{M_1}$$

            Working with the mass function,

            $$\frac{1}{f} = \frac{(1+M_1/M_2)^2}{M_2\sin^3 i} = \frac{(1+a_2/a_1)^2}{M_2\sin^3i}$$

            Substituting $(1+a_2/a_1)$

            $$\frac{1}{f} = \frac{G(M_1+M_2)}{(a_1+a_2)M_2v_1^2\sin^3i} = \frac{G(1+M_1/M_2)}{(a_1+a_2)v_1^2\sin^3i} = \frac{G(1+a_2/a_1)}{(a_1+a_2)v_1^2\sin^3i}$$

            Substituting $(1+a_2/a_1)$ again,

            $$\frac{1}{f} = \frac{G}{v_1^3\sin^3i}\sqrt{\frac{G(M_1+M_2)}{(a_1+a_2)^3}}$$

            But look at the similarity between the term inside the square root and Kepler's third law! Thus,

            $$\frac{1}{f} = \frac{2\pi G }{T v_1^3\sin^2i}$$

            Finally, substituting \(K\) and rearranging,

            $$\boxed{f = \frac{TK^3}{2\pi G}}$$

            \item Assuming $i = 90^\circ$, we have
            
            $$f = \frac{M_2^3}{(M_1+M_2)^2}$$

            Using units of solar masses and $M_1 = 2M_\odot$, we have

            $$46.2 = \frac{M_2^3}{(2+M_2)^2}$$

            $$M_2 = \left(46.2 (2+M_2)^2\right)^{1/3}$$

            Using iteration we can find $\boxed{M_2 = 49.7 M_\odot }$

            \item Since the mass of \(M_2\) is very large, it is more likely that it is a black hole.
        \end{alternativas}
    \end{pssolution*}
\end{pproblem}

\pts{3}
\begin{pproblem} Juventino is located at his (not so) secret base at the terrestrial North Pole. After an emergency call from Deduardo Letodo, he needs to send a missile to the equator in order to blow up the base of his archenemy: Raposo. 
    \begin{alternativas}
        \item Junvetino is not in his best form, and with all his strength he can only throw the rocket at the first cosmic velocity (the orbital speed for bodies near Earth's surface). Under these conditions, what is the semimajor axis of the orbit of the rocket launched by Juventino?
        \item What is the greatest height the rocket reaches?
        \item What is the time interval between Juventino launching the rocket and Raposo's base being blown up?
    \end{alternativas}

    \begin{pssolution*}{}{}
        \begin{alternativas}
            \item Using conservation of energy,
            
            \[-\frac{GM_\oplus m}{R_\oplus} + \frac{mv_0^2}{2} = -\frac{GM_\oplus m}{2a}\]

            Solving for \(a\),

            \[a = \frac{GM_\oplus}{2GM_\oplus - v_0^2 R_\oplus}R_\oplus\]

            Substituting \(v_0^2 = \frac{GM_\oplus}{R_\oplus}\), we obtain

            \[\boxed{a = R_\oplus}\]

            \item Sketching the rocket's orbit,
            
            \begin{figure}[H]
                \centering
                \includegraphics[width=0.7\linewidth]{imagens/orbitapn-eq.png}                
            \end{figure}

            By geometry, we can say that \(\overline{AO} = \overline{BO} = \frac{\overline{CD}}{2} = R_\oplus\). Using basic trigonometry, we obtain \(\overline{AB} = R_\oplus \sqrt 2\). The greatest height the rocket reaches is given by \(\overline{OD}-R_\oplus\).

            Using the property of the ellipse that the sum of the distances from a point to the two foci is \(2a\). Since \(O\) is one of the foci, the other focus, \(O'\), is also a distance \(R_\oplus\) from \(A\) and from \(B\). Under these conditions, the point \(O'\) must lie a distance \(R_\oplus\sqrt 2\) from \(O\) and on the line \(\overline{CD}\). Thus, \(AOBO'\) forms a square of side \(R_\oplus\).

            The midpoint of the segment \(\overline{OO'}\) is the center of the ellipse. Therefore, \(\overline{AB} = 2b = 2a\sqrt{1-e^2}\). Substituting the value of \(\overline{AB}\), we have
            
            \[R_\oplus\sqrt 2 = 2 R_\oplus \sqrt{1-e^2}\]

            Solving for \(e\),

            \[e = \frac{1}{\sqrt{2}}\]

            Therefore, \(\overline{OO'} = 2c = 2R_\oplus/\sqrt2\). Let \(\overline{OC} = x\). By symmetry, \(\overline{OC} = \overline{O'D}\). Setting up the equation,

            \[\overline{CD} = \overline{CO}+\overline{OO'}+\overline{O'D} = 2x + \frac{2}{\sqrt2}R_\oplus = 2R_\oplus\]

            \[x = R_\oplus\left(1 - \frac{1}{\sqrt2}\right)\]

            Now we can use \(\overline{OD} = \overline{CD} - \overline{OC} = R_\oplus\left(1+1/\sqrt2\right)\), and therefore the greatest height the rocket reaches is

            \[\boxed{h = \overline{OD}-R_\oplus = \frac{R_\oplus}{\sqrt2}}\]

            \item The point \(A\) has true anomaly \(\theta_A = \frac{\pi}{2}+\frac{\pi}{4} = \frac{3\pi}{4}\). By symmetry, the time to go from \(C\) to \(A\) is the same as the time to go from \(B\) to \(C\). Therefore, the time to go from \(A\) to \(B\) is the orbital period minus twice the time from \(O\) to \(A\). To calculate this last quantity, we will use Kepler's equation,
            
            \[M = E - e\sin E\]

            where \(E\) is the eccentric anomaly, defined by

            \[r(E)= a(1-e\cos E)\]

            \[E = \cos^{-1}\left(\frac{1-r/a}{e}\right)\]

            Substituting \(a = R_\oplus\), we can find \(E_A\). When the rocket is at the point \(A\), \(r(A)=R_\oplus\), so

            \[E_A=\frac{\pi}{2}\]

            Thus,

            \[M_A = \frac{\pi}{2}-\frac{1}{\sqrt{2}}\]

            Finally, the time to go from \(A\) to \(B\) is

            \[t_{AB} = T - 2t_{CA} = T\left(1 - \frac{M_A}{\pi}\right)\]
            \[t_{AB} = T\left(1-\frac{1}{2}+\frac{1}{\pi\sqrt2}\right)\]

            Using Kepler's third law,

            \[t_{AB} = \left(\frac{1}{2}+\frac{1}{\pi\sqrt2}\right)\sqrt{\frac{4\pi^2R_{\oplus}^3}{GM_\oplus}} \]

            \[\boxed{t_{AB} \approx 1^h01^m}\]
            
      
        \end{alternativas}
    \end{pssolution*}
    
\end{pproblem}

\pts{5}
\begin{pproblem}
    In this problem, our goal is to calculate the precession period of the Earth. To do so, we will consider the Earth as an ellipsoid, with polar radius $R_p$, equatorial radius $R_e$, rotation period about the polar axis $T = 23^h56^m4^s$, and uniformly distributed mass $M_\oplus$. For our model, consider that only the Sun and the Moon exert gravitational forces relevant to the problem.
    \begin{alternativas}
        \item Assuming that $R_p = R_\oplus$, what is the numerical value of $R_e$?
        
        \item The moment of inertia of a body can be calculated by $I = \int r^2 dm$. Given this, calculate the moment of inertia about the x-axis and about the y-axis of an ellipsoid given by the equation 
        $$\frac{x^2}{R_e^2} + \frac{y^2}{R_p^2} + \frac{z^2}{R_p^2}= 1$$
        With the moments of inertia calculated, for the Earth let us define the $x$-axis as the axis that contains the equator and the $y$-axis as the one that contains the poles. Calculate the quantity 
        $$\frac{I_p}{(I_p-I_e)}$$

        where $I_p$ is the polar moment of inertia (about the y-axis) and $I_e$ the equatorial moment of inertia (about the y-axis)
        

        \item Find an expression for the torque exerted by the Sun-Moon system on the Earth. (Hint: because of symmetry, many terms inside the integrals will vanish, and the final torque will have a component in only one direction in the coordinate system defined previously.)
        \item Finally, using the equation 
        $$\mathbf{\tau } = \Omega \times \mathbf{L} $$
        where $\mathbf{L}$ is the angular momentum and $\Omega$ is the precession frequency, find the precession period of the Earth.
    \end{alternativas}

    \begin{pssolution*}{}{}
        \begin{alternativas}
            \item Since the Earth is in hydrostatic equilibrium,
            
            \[\frac{GM_\oplus}{r_p} = \frac{GM_\oplus}{r_e}+\frac{1}{2}\omega^2 r_e^2\]
        
            Solving for \(r_e\),

            \[r_e = \frac{T}{2\pi}\sqrt{2GM_\oplus\left(\frac{1}{r_p}-\frac{1}{r_e}\right)}\]

            Iterating, we can obtain

            \[r_e \approx 6382 \text{ km}\]

            \item Given the definition of the ellipsoid, \(x\) and \(y\) lie in the equatorial plane and \(z\) passes through the center of the Earth and the north celestial pole. Calculating the moments of inertia,
            
            \[I_{zz} = \int(y^2+z^2)dm = \rho \iiint (x^2+y^2)dxdydz\]

            Using the transformation,

            \[x = R_e u \, \ y = R_e v \, \ z = R_p w\]

            so that \(u^2+v^2+w^2=1\), we have

            \[I_{zz} = \rho\iiint (R_e^2u^2+R_e^2v^2) R_e^2R_pdudvdw = \rho R_e^4R_p \iiint (u^2+v^2)dudvdw\]

            Since \(u^2+v^2+w^2=1\) describes a unit sphere, we have symmetry, which guarantees

            \[\iiint u^2d\Omega = \iint v^2d\Omega = \iint w^2d\Omega = \frac{1}{3}\iiint(u^2+v^2+w^2)d\Omega\]

            We can simplify

            \[\frac{1}{3}\iiint(u^2+v^2+w^2)d\Omega = \frac{1}{3}\int_0^{2\pi}\int_{0}^{\pi}\int_0^1 r^4\sin\theta drd\theta d\phi = \frac{4\pi}{15}\]

            Solving for \(I_{zz}\),

            \[I_{zz} = \rho R_e^4 R_p \left(\frac{4\pi}{15}+\frac{4\pi}{15}\right)\]

            Substituting \(\rho = 3M_\oplus/4\pi R_e^2R_p\), we obtain

            \[I_{zz} = \frac{2}{5}M_\oplus R_e^2\]

            Doing the same for

            \[I_{xx} = \rho R_e^2R_p\iiint ((R_ev)^2+(R_pw)^2)dudvdw\]

            we obtain

            \[I_{xx} = \frac{1}{5}M_\oplus(R_e^2+R_p^2)\]

            Also note that \(I_{zz} \equiv I_p\) and \(I_{xx}\equiv I_e\). Thus,

            \[\boxed{\frac{I_e}{I_p-I_e} \approx 580}\]

            Since our value is \(93\%\) larger than expected, in reality this quantity must correspond to \(\approx 300\). Let us call this quantity \(f\) and proceed with it from here on.

            \item Since the torque is such that
            
            \[\tau \propto \frac{M}{r^3}\]

            we have

            \[\frac{\tau_{\mathrm{Moon}}}{\tau_\odot} = \frac{M_{\mathrm{Moon}}}{M_\odot}\frac{d_\odot}{d_{\mathrm{Moon}}}\approx 2.17\]

            Now, to find the torque due to the Sun, we first calculate the force it exerts on a mass element \(dm\) in the Earth. Let us define \(\mathbf R\) as the position vector of the Sun and \(\mathbf r\) as the position vector of the mass element.

            \[d\mathbf F = -\frac{GM_\odot}{|\mathbf R - \mathbf r|^3}(\mathbf R - \mathbf r)dm\]

            With respect to the center of the Earth, the torque is

            \[d\vec\tau = \mathbf r \times d\mathbf F = \frac{GM_\odot}{|\mathbf R - \mathbf r|^3}(\mathbf R \times \mathbf r)dm\]

            Using \(R>>r\) and the approximation \((1+x)^n\approx 1+nx\), we obtain

            \[d\vec\tau \approx \frac{GM_\oplus}{R^3}\left(1+3\frac{\mathbf R \cdot \mathbf r}{R^2}\right)(\mathbf R \times \mathbf r )dm\]

            Expanding the cross product,

            \[\mathbf R \times \mathbf r = (R_yr_z - R_z r_y )\mathbf{\hat x} + (R_zr_x - R_xr_z)\mathbf {\hat y} + (R_xr_y - R_yr_x) \mathbf{\hat z}\]

            In this way,

            \[\tau_x = \frac{GM_\oplus}{R^3}\int \left(1+3\frac{R_xr_x+R_yr_y+R_zr_z}{R^2}\right)(R_yr_z-R_zr_y)dm\]

            However, when we evaluate the integral, we can see that all terms of the form \(\int r_i dm \) and \(\int r_ir_jdm\) vanish, due to symmetry. Therefore our integral reduces to

            \[\tau_x = \frac{3GM_\oplus}{R^5}R_yR_z\left(\int r_y^2dm + \int r_z^2dm\right)\]

            The remaining integrals are the moments of inertia about \(y\) and \(z\). Similarly, for \(\tau_y\) and \(\tau_z\) we obtain

            \[\tau_x = \frac{3GM_\oplus}{R^5}R_yR_z (I_{yy}-I{zz})\]
            \[\tau_y = \frac{3GM_\oplus}{R^5}R_zR_x (I_{zz}-I{xx})\]
            \[\tau_z = \frac{3GM_\oplus}{R^5}R_xR_y (I_{xx}-I{yy})\]

            We can describe \(R_i\) using elliptical coordinates.

            \[R_x = R\sin\lambda \cos\epsilon\]
            \[R_y = R\cos\lambda\]
            \[R_z = R\sin\lambda\sin\epsilon\]

            where \(\lambda\) is the ecliptic longitude of the Sun. Thus,

            \[\tau_x = \frac{3GM_\oplus}{R^3}\sin\lambda\cos\lambda\sin\epsilon (I_{yy}-I_{zz})\]
            \[\tau_y = \frac{3GM_\oplus}{R^3} \sin^2\lambda\sin\epsilon\cos\epsilon(I_{zz}-I_{xx})\]
            \[\tau_z = \frac{3GM_\oplus}{R^3} \sin\lambda\cos\lambda\cos\epsilon(I_{xx}-I_{yy})\]

            For the precession effect, what matters is the average value of the torque, but note that \[\langle\sin\lambda\cos\lambda\rangle = 0\]. Thus, the only torque that remains is that about the \(y\)-axis.

            \[\vec\tau  = \langle\tau_y\rangle\mathbf{\hat{y}} = \frac{3GM_\oplus}{2R^3}\sin\epsilon\cos\epsilon (I_{zz}-I_{xx})\mathbf{\hat{y}}\]

            This is only the contribution of the Sun. Adding that of the Moon, we have

            \[\boxed{\vec\tau_{total} = \frac{4.75GM_\oplus}{R^3}\sin\epsilon\cos\epsilon (I_{p}-I_{e})\mathbf{\hat{y}}}\]

            \item Using the given relation,
            
            \[\vec\tau_{total} = \vec\Omega \times \vec L = \Omega L_\perp\]

            The value of \(L_\perp\) is \(I_p\omega\sin\epsilon\). Thus,

            \[\frac{4.75GM_\oplus}{R^3}\sin\epsilon\cos\epsilon (I_{p}-I_{e}) = \frac{2\pi}{T}I_p\omega\sin\epsilon\]

            Solving for \(T\),

            \[\boxed{T = \frac{2\pi}{4.75}\frac{R^3\omega}{GM_\oplus\cos\epsilon}\frac{I_p}{I_p-I_e}\approx 25300\text{ years}}\]

            which is an estimate with an error of only \(2\%\) :)
        
        \end{alternativas}
    \end{pssolution*}

    \begin{pproblem}{}
        (Magna notes)
        
    \end{pproblem}
\end{pproblem}
    

\end{document}
